% ==============================================================================
% MASTER DOCUMENT SKRIPSI
% Fakultas Sains dan Teknologi - Universitas Bahaudin Mudhary Madura
% Sesuai Buku Pedoman Penulisan Skripsi FST UNIBA Madura Revisi 2026
% ==============================================================================

\documentclass[12pt, a4paper, oneside]{report}

% 1. Pemuatan Modul Konfigurasi
% ==============================================================================
% PAKET-PAKET LATEX (PACKAGES)
% ==============================================================================

% Pengaturan Margin & Ukuran Kertas
\usepackage[
    a4paper,
    left=4cm,
    top=3cm,
    right=3cm,
    bottom=3cm,
    footskip=1.5cm,
    headheight=15pt
]{geometry}

% Tipografi: Times New Roman 12pt & Justifikasi Teks Rapi
\usepackage{mathptmx}
\usepackage[utf8]{inputenc}
\usepackage[T1]{fontenc}
\usepackage{microtype}

% Spasi dan Indentasi
\usepackage{setspace}
\usepackage{indentfirst} % Paragraf pertama setelah heading tetap menjorok masuk
\usepackage{enumitem}    % Pengendalian presisi spasi daftar (list/enumerate/itemize) tanpa space after/before

% Simbol & Rumus Matematika
\usepackage{amsmath}
\usepackage{amssymb}
\usepackage{amsfonts}

% Gambar dan Grafik
\usepackage{graphicx}
\usepackage{float}
\usepackage{tikz}
\usetikzlibrary{shapes.geometric, arrows.meta, positioning, fit, calc}

% Tabel Akademik
\usepackage{booktabs}
\usepackage{tabularx}
\usepackage{multirow}
\usepackage{array}
\usepackage{colortbl}

% Format Judul Bab & Subbab
\usepackage{titlesec}
\usepackage{etoolbox}

% Header dan Footer (Penomoran Halaman)
\usepackage{fancyhdr}

% Keterangan Gambar & Tabel
\usepackage{caption}
\usepackage{subcaption}

% Penyorotan Kode Program (Syntax Highlighting)
\usepackage{listings}
\usepackage{xcolor}

% Pengaturan Daftar Isi, Gambar, dan Tabel
\usepackage{titletoc}

% Sitasi & Daftar Pustaka Gaya APA (American Psychological Association)
\usepackage[natbibapa]{apacite}

% Hyperlink dokumen PDF
\usepackage{hyperref}
\usepackage{url}

% ==============================================================================
% METADATA SKRIPSI
% Fakultas Sains dan Teknologi - Universitas Bahaudin Mudhary Madura
% ==============================================================================

% Judul Skripsi (Bahasa Indonesia & Bahasa Inggris)
% Istilah bahasa asing dimiringkan (\textit{...}) agar taat kaidah bahasa baku
\newcommand{\skripsiJudul}{ANALISIS PERBANDINGAN KINERJA NGINX DAN HAPROXY SEBAGAI \textit{LOAD BALANCER} PADA \textit{WEB SERVICE} BERBASIS \textit{DOCKER} DENGAN KAPASITAS \textit{BACKEND} HETEROGEN}
\newcommand{\skripsiJudulEn}{COMPARATIVE PERFORMANCE ANALYSIS OF NGINX AND HAPROXY AS LOAD BALANCERS IN DOCKER-BASED WEB SERVICES WITH HETEROGENEOUS BACKEND CAPACITIES}

% Tipe Dokumen (PROPOSAL SKRIPSI untuk pra-eksperimen / SKRIPSI untuk tugas akhir)
\newcommand{\skripsiTipeDokumen}{PROPOSAL SKRIPSI}
\newcommand{\skripsiBidangMinat}{Jaringan Komputer dan Sistem Terdistribusi}
\newcommand{\skripsiTujuanDokumen}{Diajukan untuk memenuhi persyaratan\\ Memperoleh gelar Sarjana Komputer}

% Identitas Penulis
\newcommand{\skripsiPenulis}{Ahmad Muqtafi}
\newcommand{\skripsiNIM}{2302310264}

% Informasi Program Studi & Fakultas
\newcommand{\skripsiGelar}{Sarjana Komputer}
\newcommand{\skripsiSingkatanGelar}{S.Kom.}
\newcommand{\skripsiProdi}{Informatika}
\newcommand{\skripsiFakultas}{Fakultas Sains dan Teknologi}
\newcommand{\skripsiUniversitas}{Universitas K.H. Bahaudin Mudhary Madura}
\newcommand{\skripsiKota}{Sumenep}
\newcommand{\skripsiTahun}{2026}
\newcommand{\skripsiBulanTahun}{Mei 2026}

% Dosen Pembimbing
\newcommand{\skripsiPembimbingSatu}{Prof. Dr. Ir. Rachmad Hidayat, M.T., IPU., ASEAN Eng.}
\newcommand{\skripsiNIKPembimbingSatu}{19740619.201808.1.010}
\newcommand{\skripsiPembimbingDua}{Akhmad Tajuddin Tholaby M.S., S.Kom., M.Kom.}
\newcommand{\skripsiNIKPembimbingDua}{19890330.202203.1.079}

% Pimpinan Fakultas & Program Studi
\newcommand{\skripsiDekan}{Emon Rifa'i, S.T., M.T.}
\newcommand{\skripsiKaprodi}{Zeinor Rahman, S.Pd., M.Pd.}
\newcommand{\skripsiNIKKaprodi}{19960306.202108.1.061}

% Tanggal Pengesahan (Format Bahasa Indonesia)
\newcommand{\tanggalIndonesia}{%
  \the\day\ \ifcase\month\or Januari\or Februari\or Maret\or April\or Mei\or Juni\or Juli\or Agustus\or September\or Oktober\or November\or Desember\fi\ \the\year%
}
\newcommand{\skripsiTanggalPengesahan}{\tanggalIndonesia}

% ==============================================================================
% PENGATURAN FORMATTING & TATA LETAK SKRIPSI
% Sesuai Pedoman Penulisan Skripsi FST UNIBA Madura 2026
% ==============================================================================

% 1. SPASI & INDENTASI
\onehalfspacing                      % Spasi antar baris 1.5
\setlength{\parindent}{0.75cm}       % Indent alinea baru 0.75 cm (7 ketukan)
\setlength{\parskip}{0pt}            % Jarak antar paragraf 0pt (tanpa spasi kosong ekstra antar paragraf)
\raggedbottom                        % Cegah peregangan elastis spasi vertikal LaTeX (persis perilaku Word DOCX)

% Pengaturan Spasi Daftar (List: Enumerate & Itemize) tanpa spasi ekstra (identik format Word DOCX)
\setlist{
  nosep,                             % Hapus space before/after dan space ekstra antar butir
  topsep=0pt,
  partopsep=0pt,
  parsep=0pt,
  itemsep=0pt,
  leftmargin=0.75cm,                 % Indentasi sejajar alinea baru
  labelsep=0.25cm
}
\setlist[enumerate]{label=\arabic*., ref=\arabic*}
\setlist[itemize]{label=\textbullet}

% 2. PENAMAAN STRUKTUR DOKUMEN (BAHASA INDONESIA)
\renewcommand{\contentsname}{DAFTAR ISI}
\renewcommand{\listfigurename}{DAFTAR GAMBAR}
\renewcommand{\listtablename}{DAFTAR TABEL}
\renewcommand{\bibname}{DAFTAR PUSTAKA}
\renewcommand{\chaptername}{BAB}
\renewcommand{\appendixname}{LAMPIRAN}
\renewcommand{\figurename}{Gambar}
\renewcommand{\tablename}{Tabel}

% 3. FORMAT PENOMORAN PERSAMAAN, TABEL, DAN GAMBAR PER BAB
% Sesuai pedoman: Persamaan (3-18), Tabel 3.1, Gambar 3.1
\renewcommand{\thechapter}{\Roman{chapter}}
\renewcommand{\thesection}{\arabic{chapter}.\arabic{section}}
\renewcommand{\thesubsection}{\arabic{chapter}.\arabic{section}.\arabic{subsection}}
\renewcommand{\theequation}{\arabic{chapter}-\arabic{equation}}
\renewcommand{\thefigure}{\arabic{chapter}.\arabic{figure}}
\renewcommand{\thetable}{\arabic{chapter}.\arabic{table}}

% 4. FORMAT JUDUL BAB, SUB-BAB, DAN SUB-SUB-BAB (TITLESEC)
% Format Bab: BAB I <baris baru> JUDUL BAB (Kapital, Tebal, Rata Tengah; BAB X: 14pt, Judul: 12pt)
\titleformat{\chapter}[display]
  {\singlespacing\centering\bfseries}
  {\fontsize{14pt}{18pt}\selectfont\MakeUppercase{\chaptername\ \thechapter}}
  {0.5em}
  {\fontsize{12pt}{16pt}\selectfont\MakeUppercase}

% Spasi bab: tepat 3.0 cm dari sisi atas kertas (persis format Word proposal teman)
\titlespacing*{\chapter}{0pt}{-35pt}{1.5em}

% Format Sub-bab (Section): Tebal, Rata Kiri, Title Case (12pt)
\titleformat{\section}
  {\singlespacing\bfseries\fontsize{12pt}{16pt}\selectfont}
  {\thesection}
  {0.75em}
  {}
\titlespacing*{\section}{0pt}{1.0em}{0.2em}

% Format Sub-sub-bab (Subsection): Tebal, Rata Kiri (12pt)
\titleformat{\subsection}
  {\singlespacing\bfseries\fontsize{12pt}{16pt}\selectfont}
  {\thesubsection}
  {0.75em}
  {}
\titlespacing*{\subsection}{0pt}{0.8em}{0.2em}

% Otomatisasi eliminasi jeda spasi ekstra jika \subsection langsung mengikuti \section tanpa teks perantara
% (menghilangkan kesan adanya baris kosong atau space after paragraph antara 2.2 dengan 2.2.1 / 2.3 dengan 2.3.1)
\newif\ifdirectlyaftersection
\directlyaftersectionfalse

\pretocmd{\section}{\directlyaftersectiontrue}{}{}
\pretocmd{\subsection}{%
  \ifdirectlyaftersection
    \vspace{-0.85em}%
    \directlyaftersectionfalse
  \fi
}{}{}

\makeatletter
\everypar=\expandafter{\the\everypar \directlyaftersectionfalse}
\makeatother

% 5. LOGIKA PENOMORAN HALAMAN (FANCYHDR)
% Pedoman FST UNIBA 2026:
% - Bagian awal: Angka Romawi kecil di tengah bawah
% - Bagian utama: Angka Arab di kanan atas
% - Halaman pertama setiap bab baru: Angka Arab di tengah bawah

% Style halaman awal skripsi (Frontmatter)
\fancypagestyle{frontstyle}{
  \fancyhf{}
  \renewcommand{\headrulewidth}{0pt}
  \renewcommand{\footrulewidth}{0pt}
  \fancyfoot[C]{\thepage}
}

% Style halaman isi utama (Mainmatter)
\fancypagestyle{mainstyle}{
  \fancyhf{}
  \renewcommand{\headrulewidth}{0pt}
  \renewcommand{\footrulewidth}{0pt}
  \fancyhead[R]{\thepage}
}

% Style 'plain' otomatis dipanggil pada halaman pertama \chapter
\fancypagestyle{plain}{
  \fancyhf{}
  \renewcommand{\headrulewidth}{0pt}
  \renewcommand{\footrulewidth}{0pt}
  \fancyfoot[C]{\thepage}
}

% 6. PENGATURAN SPASI & CAPTION OBJEK GAMBAR & TABEL
% Pengendalian spasi sebelum dan sesudah float/in-text float (menghilangkan jeda ekstra seperti space after)
\setlength{\floatsep}{4pt plus 1pt minus 1pt}
\setlength{\textfloatsep}{6pt plus 1pt minus 1pt}
\setlength{\intextsep}{0pt}
\setlength{\dblfloatsep}{4pt plus 1pt minus 1pt}
\setlength{\dbltextfloatsep}{6pt plus 1pt minus 1pt}

% Tabel: Judul di atas tabel, jarak dengan tabel rapat, tanpa spasi ekstra di atas judul
\captionsetup[table]{
  labelfont={bf},
  textfont={normalfont},
  justification=centering,
  labelsep=space,
  position=top,
  aboveskip=0pt,
  belowskip=4pt
}

% Gambar: Judul di bawah gambar, huruf label tebal, tanpa spasi ekstra di bawah judul
\captionsetup[figure]{
  labelfont={bf},
  textfont={normalfont},
  justification=centering,
  labelsep=space,
  position=bottom,
  aboveskip=4pt,
  belowskip=0pt
}

% 7. FORMAT KUTIPAN DAFTAR PUSTAKA APA
\AtBeginDocument{
  \renewcommand{\bibname}{DAFTAR PUSTAKA}
  \renewcommand{\refname}{DAFTAR PUSTAKA}
}
\setlength{\bibhang}{0.75cm} % Hanging indent 0.75 cm
\setlength{\bibsep}{6.9pt}   % Jarak jeda antar-entri pustaka (identik dengan after=138 twips / 6.9 pt pada acuan Word)

% 8. PENYOROTAN KODE PROGRAM (LISTINGS)
\definecolor{codegray}{rgb}{0.5,0.5,0.5}
\definecolor{backcolour}{rgb}{0.97,0.97,0.97}
\definecolor{primaryblue}{rgb}{0.0,0.2,0.6}
\definecolor{stringgreen}{rgb}{0.1,0.5,0.2}

\lstdefinestyle{skripsicode}{
  backgroundcolor=\color{backcolour},
  commentstyle=\color{codegray}\itshape,
  keywordstyle=\color{primaryblue}\bfseries,
  numberstyle=\tiny\color{codegray},
  stringstyle=\color{stringgreen},
  basicstyle=\ttfamily\footnotesize\singlespacing,
  breakatwhitespace=false,
  breaklines=true,
  captionpos=b,
  keepspaces=true,
  numbers=left,
  numbersep=7pt,
  showspaces=false,
  showstringspaces=false,
  showtabs=false,
  tabsize=2,
  frame=single,
  rulecolor=\color{lightgray}
}
\lstset{style=skripsicode}

% 9. HYPERLINK SETUP (PDF METADATA & LINK COLORING)
\hypersetup{
  colorlinks=true,
  linkcolor=black,
  citecolor=black,
  urlcolor=blue,
  pdfborder={0 0 0}
}

% 10. FORMAT DAFTAR ISI, DAFTAR GAMBAR, DAN DAFTAR TABEL (TITLETOC)
% Format Titik-Titik Rapat Bersambung Penuh Hingga Nomor Halaman & Hanging Indent
\contentsmargin[1.5em]{1.2em}

% Format Bab / Bagian Utama (BAB I PENDAHULUAN ...................... 1)
% Untuk entri tanpa nomor (frontmatter, daftar pustaka, lampiran): tercetak biasa tebal dengan titik rapat
\titlecontents{chapter}[0pt]
  {\vspace{0.35em}\bfseries}
  {BAB \thecontentslabel\ }
  {}
  {\titlerule*[0.28pc]{.}\contentspage}

% Format Sub-bab / Section (1.1 Latar Belakang ..................... 1)
% Indent label 0.8cm, teks judul & baris lipat mulai pada 1.6cm (hanging indent)
\titlecontents{section}[1.6cm]
  {\normalfont}
  {\contentslabel[\thecontentslabel]{0.8cm}}
  {}
  {\titlerule*[0.28pc]{.}\contentspage}

% Format Sub-sub-bab / Subsection (2.2.1 Sistem .................... 6)
% Indent label 1.6cm, teks judul & baris lipat mulai pada 2.8cm (hanging indent)
\titlecontents{subsection}[2.8cm]
  {\normalfont}
  {\contentslabel[\thecontentslabel]{1.2cm}}
  {}
  {\titlerule*[0.28pc]{.}\contentspage}

% Format Gambar pada Daftar Gambar (Gambar 3.1 Topologi ........... 11)
% Label "Gambar X.X" lebar 2.5cm, baris lipat hanging indent sejajar pada 2.5cm
\titlecontents{figure}[2.5cm]
  {\normalfont}
  {\contentslabel[\textbf{Gambar \thecontentslabel}]{2.5cm}}
  {}
  {\titlerule*[0.28pc]{.}\contentspage}

% Format Tabel pada Daftar Tabel (Tabel 2.1 Matriks ................ 6)
% Label "Tabel X.X" lebar 2.2cm, baris lipat hanging indent sejajar pada 2.2cm
\titlecontents{table}[2.2cm]
  {\normalfont}
  {\contentslabel[\textbf{Tabel \thecontentslabel}]{2.2cm}}
  {}
  {\titlerule*[0.28pc]{.}\contentspage}




\begin{document}

% ==============================================================================
% BAGIAN AWAL (FRONTMATTER)
% ==============================================================================
\pagestyle{empty}
\hypersetup{pageanchor=false}

% Halaman Sampul (Cover Proposal) - Tanpa Nomor Halaman (Halaman i tersirat)
% ==============================================================================
% HALAMAN SAMPUL (COVER)
% Sesuai Lampiran 1b Pedoman Penulisan Skripsi FST UNIBA Madura 2026
% ==============================================================================

\begin{titlepage}
\begin{center}
    \singlespacing
    
    % Judul Skripsi (Kapital, Tebal, 12pt proporsional)
    {\fontsize{12pt}{16pt}\selectfont \bfseries \skripsiJudul}
    
    \vspace{1.5cm}
    
    % Tulisan Tipe Dokumen (PROPOSAL SKRIPSI / SKRIPSI) & Bidang Minat
    {\fontsize{12pt}{16pt}\selectfont \bfseries \MakeUppercase{\skripsiTipeDokumen}}\\[0.25cm]
    {\fontsize{12pt}{16pt}\selectfont \skripsiBidangMinat}
    
    \vspace{0.9cm}
    
    % Kalimat Persyaratan (12pt reguler)
    {\fontsize{12pt}{16pt}\selectfont
    \skripsiTujuanDokumen
    }
    
    \vspace{1.2cm}
    
    % Logo Universitas Bahaudin Mudhary Madura
    \includegraphics[width=4.0cm]{figures/logo-uniba.png}
    
    \vspace{1.2cm}
    
    % Identitas Mahasiswa (Nama tebal, NIM reguler)
    {\fontsize{12pt}{16pt}\selectfont \bfseries \MakeUppercase{\skripsiPenulis}}\\
    \vspace{0.2cm}
    {\fontsize{12pt}{16pt}\selectfont NIM. \skripsiNIM}
    
    \vfill
    
    % Identitas Institusi
    {\fontsize{12pt}{16pt}\selectfont \bfseries
    PROGRAM STUDI \MakeUppercase{\skripsiProdi}\\
    \MakeUppercase{\skripsiFakultas}\\
    \MakeUppercase{\skripsiUniversitas}\\
    \MakeUppercase{\skripsiKota}\\
    \skripsiTahun
    }

\end{center}
\end{titlepage}


% Aktifkan pageanchor untuk hyperlink dokumen PDF
\hypersetup{pageanchor=true}

% Lembar Pengesahan Dosen Pembimbing (Halaman ii)
% ==============================================================================
% LEMBAR PENGESAHAN DOSEN PEMBIMBING
% Sesuai Format Skripsi Kating (FST UNIBA Madura)
% ==============================================================================

\clearpage
\phantomsection
\pagenumbering{roman}
\setcounter{page}{2}
\pagestyle{frontstyle}
\addcontentsline{toc}{chapter}{LEMBAR PENGESAHAN}

\begin{center}
    \singlespacing
    
    {\fontsize{12pt}{16pt}\selectfont \bfseries LEMBAR PENGESAHAN}
    
    \vspace{0.65cm}
    
    {\fontsize{12pt}{16pt}\selectfont \bfseries \skripsiJudul}
    
    \vspace{0.6cm}
    
    {\fontsize{12pt}{16pt}\selectfont \bfseries \MakeUppercase{\skripsiTipeDokumen}}\\[0.15cm]
    {\fontsize{12pt}{16pt}\selectfont \skripsiBidangMinat}
    
    \vspace{0.55cm}
    
    {\fontsize{12pt}{16pt}\selectfont
    \skripsiTujuanDokumen
    }
    
    \vspace{0.6cm}
    
    % Logo Universitas Bahaudin Mudhary Madura (Diperbesar agar proporsional)
    \includegraphics[width=3.5cm]{figures/logo-uniba.png}
    
    \vspace{0.6cm}
    
    {\fontsize{12pt}{16pt}\selectfont \MakeUppercase{\skripsiPenulis}}\\[0.12cm]
    {\fontsize{12pt}{16pt}\selectfont NIM. \skripsiNIM}
    
    \vspace{0.6cm}
    
    {\fontsize{12pt}{16pt}\selectfont
    Proposal Skripsi ini telah disetujui oleh dosen pembimbing\\
    pada tanggal \dots\ \dots\ \skripsiTahun
    }
    
    \vspace{0.95cm}
    
    % Kolom Pembimbing I dan II (Ruang tanda tangan 1.6cm)
    \noindent
    \begin{minipage}[t]{0.48\textwidth}
        \centering
        Dosen Pembimbing I\\[1.6cm]
        \underline{Prof. Dr. Ir. Rachmad Hidayat, M.T.,}\\
        \underline{IPU., ASEAN Eng.}\\[2pt]
        NIP. \skripsiNIKPembimbingSatu
    \end{minipage}\hfill
    \begin{minipage}[t]{0.48\textwidth}
        \centering
        Dosen Pembimbing II\\[1.6cm]
        \underline{Akhmad Tajuddin Tholaby M.S.,}\\
        \underline{S.Kom., M.Kom.}\\[2pt]
        NIP. \skripsiNIKPembimbingDua
    \end{minipage}
    
    \vspace{0.95cm}
    
    % Mengetahui Kaprodi (Ruang tanda tangan 1.6cm)
    Mengetahui,\\
    Ketua Program Studi\\[1.6cm]
    \underline{\skripsiKaprodi}\\[2pt]
    NIP. \skripsiNIKKaprodi

\end{center}


% Bagian-bagian yang hanya untuk Naskah Skripsi Lengkap Pasca-Eksperimen
% dinonaktifkan sementara untuk tahap Proposal Skripsi (Pra-Eksperimen / Bab 1-3).
% Buka tanda komentar (%) di bawah ini saat menyusun naskah Skripsi akhir:
% % ==============================================================================
% PERNYATAAN ORISINILITAS SKRIPSI
% Sesuai Format Skripsi Kating (FST UNIBA Madura)
% ==============================================================================

\clearpage
\phantomsection
\pagestyle{frontstyle}
\addcontentsline{toc}{chapter}{PERNYATAAN ORISINILITAS SKRIPSI}

\begin{center}
    {\fontsize{14pt}{18pt}\selectfont \bfseries PERNYATAAN ORISINILITAS SKRIPSI}
\end{center}

\vspace{1.0cm}

Saya menyatakan dengan sebenar-benarnya bahwa sepanjang pengetahuan saya dan berdasarkan hasil penelusuran berbagai karya ilmiah, gagasan dan masalah ilmiah yang diteliti dan diulas di dalam Naskah Skripsi ini adalah asli dari pemikiran saya. tidak terdapat karya ilmiah yang pernah diajukan oleh orang lain untuk memperoleh gelar akademik di suatu Perguruan Tinggi, dan tidak terdapat karya atau pendapat yang pernah ditulis atau diterbitkan oleh orang lain, kecuali yang secara tertulis dikutip dalam naskah ini dan disebutkan dalam sumber kutipan dan daftar pustaka.

Apabila ternyata di dalam naskah Skripsi/Disertasi ini dapat dibuktikan terdapat unsur-unsur jiplakan, saya bersedia Skripsi dibatalkan, serta diproses sesuai dengan peraturan perundang-undangan yang berlaku (UU No. 20 Tahun 2003, pasal 25 ayat 2 dan pasal 70).

\vspace{2.0cm}

\begin{flushright}
\begin{tabular}{l}
    \skripsiKota, \dots\ \dots\ \skripsiTahun\\
    Mahasiswa,\\[2.2cm]
    \skripsiPenulis\\
    NIM. \skripsiNIM
\end{tabular}
\end{flushright}

% % ==============================================================================
% ABSTRAK (BAHASA INDONESIA)
% Sesuai Format Skripsi Kating (FST UNIBA Madura)
% ==============================================================================

\clearpage
\phantomsection
\pagestyle{frontstyle}
\addcontentsline{toc}{chapter}{ABSTRAK}

\begin{center}
    {\fontsize{14pt}{18pt}\selectfont \bfseries ABSTRAK}
\end{center}

\vspace{0.8cm}

\noindent
\skripsiPenulis, Program Studi \skripsiProdi, \skripsiFakultas, \skripsiUniversitas, \dots\ \skripsiTahun, \textit{\skripsiJudul}.\\
Dosen Pembimbing: \skripsiPembimbingSatu\ dan \skripsiPembimbingDua.

\vspace{0.6cm}

\begin{singlespace}
Perkembangan teknologi komputasi awan dan arsitektur \textit{microservices} telah mendorong pemanfaatan kontainerisasi berbasis Docker sebagai tulang punggung penyedia layanan \textit{web service} modern. Dalam operasional lingkungan produksi sesungguhnya, kluster kontainer sering kali berada dalam kondisi kapasitas sumber daya yang heterogen akibat pembatasan kuota \textit{Control Groups} (cgroups) kernel Linux maupun ketimpangan spesifikasi perangkat keras. Kondisi asimetris ini menimbulkan tantangan krusial pada penyeimbang beban (\textit{load balancer}), khususnya ketika melayani karakteristik beban kerja intensif prosesor (\textit{CPU-bound}) dibandingkan beban kerja intensif masukan/keluaran (\textit{I/O-bound}). Penelitian ini bertujuan untuk menganalisis perbandingan kinerja data-plane NGINX dan HAProxy menggunakan empat matriks algoritma penjadwalan distribusi beban: \textit{Round Robin}, \textit{Least Connections}, \textit{Weighted Round Robin}, dan \textit{Weighted Least Connections} pada kluster backend heterogen. Parameter evaluasi yang diukur mencakup \textit{throughput} (req/s), rerata latensi respon, \textit{tail latency} persentil p95 dan p99, persentase kegagalan koneksi (\textit{error rate}), serta utilisasi vCPU dan RAM. Hasil benchmarking empiris menunjukkan bahwa kesesuaian antara penentuan bobot algoritma (\textit{weighted}) dengan alokasi kuota cgroups kontainer secara signifikan mampu mengeliminasi fenomena \textit{CPU throttling} dan menekan lonjakan latensi ekor pada beban konkurensi tinggi.
\end{singlespace}

\vspace{0.8cm}

\noindent
\textbf{Kata Kunci:} NGINX, HAProxy, Load Balancer, Docker, Backend Heterogen, Algoritma Penjadwalan

% % ==============================================================================
% ABSTRACT (BAHASA INGGRIS)
% Sesuai Format Skripsi Kating (FST UNIBA Madura)
% ==============================================================================

\clearpage
\phantomsection
\pagestyle{frontstyle}
\addcontentsline{toc}{chapter}{\textit{ABSTRACT}}

\begin{center}
    {\fontsize{14pt}{18pt}\selectfont \bfseries ABSTRACT}
\end{center}

\vspace{0.8cm}

\noindent
\skripsiPenulis, Informatics Study Program, Faculty of Science and Technology, \skripsiUniversitas, \dots\ \skripsiTahun, \textit{\skripsiJudulEn}.\\
Advisors: \skripsiPembimbingSatu\ and \skripsiPembimbingDua.

\vspace{0.6cm}

\begin{singlespace}
\textit{The advancement of cloud computing and microservices architecture has driven the widespread adoption of Docker-based containerization as the backbone of modern web services. In real-world production environments, backend container clusters often operate under heterogeneous resource capacities due to Linux kernel Control Groups (cgroups) quota allocations or hardware discrepancies. This asymmetric condition introduces critical challenges to load balancers, particularly when handling processor-intensive (CPU-bound) workloads compared to input/output-intensive (I/O-bound) traffic. This study aims to analyze the comparative performance of NGINX and HAProxy data-plane engines across four scheduling algorithm matrices: Round Robin, Least Connections, Weighted Round Robin, and Weighted Least Connections within a heterogeneous backend container cluster. Evaluated performance metrics encompass throughput (requests per second), average response latency, 95th and 99th percentile tail latency, connection error rates, as well as vCPU and memory utilization. The empirical benchmarking results demonstrate that aligning weighted scheduling algorithms with container cgroups quotas significantly mitigates CPU throttling phenomena and prevents tail latency degradation under high concurrency stress.}
\end{singlespace}

\vspace{0.8cm}

\noindent
\textit{\textbf{Keywords:} NGINX, HAProxy, Load Balancer, Docker, Heterogeneous Backend, Scheduling Algorithms}


% Kata Pengantar (Halaman iii)
% ==============================================================================
% KATA PENGANTAR
% Sesuai Format Skripsi Kating (FST UNIBA Madura)
% ==============================================================================

\clearpage
\phantomsection
\pagestyle{frontstyle}
\addcontentsline{toc}{chapter}{KATA PENGANTAR}

\begin{center}
    {\fontsize{14pt}{18pt}\selectfont \bfseries KATA PENGANTAR}
\end{center}

\vspace{0.3cm}

\textit{Alhamdulillahi Rabbil 'Alamin}, segala puji dan syukur penulis panjatkan kehadirat Allah SWT atas limpahan berkah, rahmat, dan karunia-Nya, sehingga penulis dapat menyelesaikan penulisan proposal skripsi yang berjudul \textbf{“\skripsiJudul”} ini sebagai salah satu syarat untuk memperoleh gelar \skripsiGelar\ (\skripsiSingkatanGelar) pada Program Studi \skripsiProdi, \skripsiFakultas, \skripsiUniversitas.

Penyusunan skripsi ini tidak lepas dari bimbingan, arahan, dan dukungan dari berbagai pihak. Oleh karena itu, pada kesempatan ini penulis menyampaikan terima kasih yang tulus kepada:

\begin{enumerate}\singlespacing\setlength{\itemsep}{1.5pt}\setlength{\parsep}{0pt}\setlength{\topsep}{2pt}
    \item Bapak \skripsiDekan, selaku Dekan Fakultas Sains dan Teknologi Universitas K.H. Bahaudin Mudhary Madura.
    \item Bapak Mohammad Iqbal Bachtiar, S.T., M.T., selaku Wakil Dekan Fakultas Sains dan Teknologi Universitas K.H. Bahaudin Mudhary Madura.
    \item Bapak \skripsiKaprodi, selaku Ketua Program Studi \skripsiProdi\ Universitas K.H. Bahaudin Mudhary Madura.
    \item Bapak \skripsiPembimbingSatu, selaku Dosen Pembimbing I atas segala bimbingan, telaah ilmiah, dan arahannya.
    \item Bapak \skripsiPembimbingDua, selaku Dosen Pembimbing II atas bimbingan teknis, saran konstruktif, dan motivasinya.
    \item Seluruh Dosen dan Tenaga Kependidikan Program Studi \skripsiProdi\ Fakultas Sains dan Teknologi.
    \item Teruntuk Kedua Orang Tua tercinta, Ayahanda dan Ibunda, serta segenap keluarga besar atas doa tulus dan pengorbanan tanpa pamrih.
    \item Rekan-rekan mahasiswa seperjuangan Program Studi \skripsiProdi\ atas kerja sama dan kebersamaannya.
\end{enumerate}

Penulis menyadari bahwa dalam penulisan naskah ini masih terdapat keterbatasan. Kritik dan saran membangun sangat penulis harapkan demi penyempurnaan karya ilmiah ini. Semoga skripsi ini bermanfaat bagi kita semua. Aamiin.

\vspace{0.3cm}

\begin{flushright}
\begin{tabular}{l}
    \skripsiKota, \dots\ \dots\ \skripsiTahun\\[1.4cm]
    \skripsiPenulis
\end{tabular}
\end{flushright}


% Daftar Isi
\clearpage
\phantomsection
\addcontentsline{toc}{chapter}{DAFTAR ISI}
{\singlespacing
\tableofcontents}

% Daftar Tabel (Sesuai Pedoman FST 2026 & Proposal Sempro ACC: Daftar Tabel mendahului Daftar Gambar)
\clearpage
\phantomsection
\addcontentsline{toc}{chapter}{DAFTAR TABEL}
{\singlespacing
\listoftables}

% Daftar Gambar
\clearpage
\phantomsection
\addcontentsline{toc}{chapter}{DAFTAR GAMBAR}
{\singlespacing
\listoffigures}

% ==============================================================================
% BAGIAN UTAMA (MAINMATTER / BAB I - III: PROPOSAL SKRIPSI)
% Penomoran angka Arab (1, 2, 3, ...) di kanan atas,
% kecuali halaman awal bab di tengah bawah
% ==============================================================================
\clearpage
\pagenumbering{arabic}
\setcounter{page}{1}
\pagestyle{mainstyle}

% Bab I s.d. Bab III (Fokus saat ini sebelum eksperimen)
% ==============================================================================
% BAB I: PENDAHULUAN
% ==============================================================================

\chapter{PENDAHULUAN}

\section{Latar Belakang}
Evolusi arsitektur \textit{web service} modern telah bergeser secara masif dari peladen fisik (\textit{bare-metal}) dan mesin virtual (\textit{Virtual Machine} / VM) menuju ekosistem kontainerisasi berbasis Docker \citep{khan2026decomposing, gurramkonda2026docker, nurzaman2022analisis}. Kontainerisasi menawarkan portabilitas tinggi, efisiensi konsumsi sumber daya komputasi, serta kecepatan \textit{deployment} aplikasi berskala besar. Namun, dalam lingkungan produksi yang sesungguhnya, kluster kontainer \textit{backend} jarang beroperasi dalam kondisi homogen yang ideal. Pembatasan kuota komputasi melalui fitur \textit{Control Groups} (cgroups) pada kernel Linux, fragmentasi perangkat keras, serta alokasi dinamis pada arsitektur \textit{microservices} menciptakan asimetri kapasitas komputasi (baik alokasi vCPU maupun kuota RAM) antar kontainer \textit{backend} \citep{shetty2025cpu, faridah2026performance}. Dalam praktik industri komputasi awan, variasi kapasitas ini umumnya mengadopsi pola tingkatan instansi berjenjang (\textit{instance tiering} atau \textit{t-shirt sizing}: \textit{Small}, \textit{Medium}, \textit{Large}) dengan rasio penggandaan kapasitas eksponensial seperti 1:2:4, yang mencerminkan heterogenitas terencana antar node layanan.

Penerapan penyeimbang beban (\textit{load balancer}) seperti NGINX dan HAProxy telah menjadi standar industri dalam mendistribusikan lalu lintas HTTP menuju jajaran peladen aplikasi \citep{buke2026load, sakib2025dynamic, arini2025perbandingan}. Kendati demikian, kedua perangkat lunak penyeimbang beban tersebut dibangun di atas filosofi arsitektur internal yang berbeda secara fundamental. NGINX mengadopsi model arsitektur \textit{multi-process} berbasis \textit{master-worker}, di mana sekumpulan proses pekerja (\textit{worker}) mengeksekusi \textit{non-blocking event-loop} secara mandiri dalam ruang memori terpisah. Sebaliknya, HAProxy dirancang murni sebagai proksi lapis aplikasi berkinerja tinggi yang beroperasi dengan model proses tunggal berorientasi peristiwa (\textit{event-driven single process}) dan memanfaatkan struktur data pohon biner elastis (\textit{elastic binary tree} / \texttt{ebtree}) berkecepatan $O(1)$ untuk manajemen antrean koneksi serta penjadwalan waktu respons deterministik \citep{tarreau2024haproxy}. Perbedaan arsitektural konkurensi ini diprediksi menghasilkan respons stabilitas yang sangat berlainan ketika penyeimbang beban berhadapan dengan kluster kontainer \textit{backend} heterogen, terlebih saat menerima karakteristik beban kerja yang bertolak belakang, yaitu beban intensif komputasi (\textit{CPU-bound}) versus beban intensif masukan/keluaran (\textit{I/O-bound}) \citep{baeturohman2024analisis}.

Mayoritas penelitian terdahulu yang membandingkan performa NGINX dan HAProxy pada lingkungan kontainer Docker hanya mengevaluasi kondisi \textit{backend} yang homogen serta terpaku pada algoritma statis tanpa bobot seperti \textit{Round Robin} \citep{nurzaman2022analisis, ramadhani2026analysis}. Di sisi lain, penelitian yang menguji algoritma berbobot (\textit{Weighted Round Robin} dan \textit{Weighted Least Connections}) pada lingkungan heterogen umumnya diimplementasikan pada Virtual Private Server (VPS) atau mesin virtual terisolasi tanpa mempertimbangkan karakteristik isolasi kernel cgroups Docker \citep{rizqi2022analisis}. Ketidaksesuaian antara logika algoritma distribusi beban dan batasan kuota cgroups berpotensi memicu fenomena \textit{CPU throttling} akibat pembatasan Completely Fair Scheduler (CFS) kernel Linux \citep{shetty2025cpu}, penumpukan antrean koneksi (\textit{backlog}), dan lonjakan ekstrem pada latensi ekor (\textit{tail latency} persentil p95 dan p99) \citep{wodislawski2026ponytail, li2025tailbench}. Oleh karena itu, diperlukan penelitian komparatif yang menguji secara empiris dan terukur kinerja data-plane NGINX versus HAProxy pada lingkungan Docker dengan kapasitas heterogen.

\section{Rumusan Masalah}
Berdasarkan latar belakang masalah yang telah diuraikan, rumusan masalah dalam penelitian ini adalah sebagai berikut:
\begin{enumerate}
    \item Bagaimana perbandingan kinerja data-plane NGINX dan HAProxy menggunakan algoritma \textit{Round Robin}, \textit{Least Connections}, \textit{Weighted Round Robin}, dan \textit{Weighted Least Connections} pada kluster kontainer Docker dengan kapasitas \textit{backend} heterogen?
    \item Bagaimana pengaruh karakteristik beban kerja intensif prosesor (\textit{CPU-bound}) dan beban kerja intensif masukan/keluaran (\textit{I/O-bound}) terhadap \textit{throughput}, latensi rata-rata, dan \textit{tail latency} (p95 dan p99) pada masing-masing mesin penyeimbang beban?
    \item Sejauh mana presisi distribusi beban aktual yang dihasilkan oleh algoritma penyeimbang beban terhadap proporsi kapasitas kuota cgroups kontainer \textit{backend}?
\end{enumerate}

\section{Batasan Masalah}
Agar penelitian ini terarah dan mencapai sasaran yang diharapkan, batasan masalah ditetapkan sebagai berikut:
\begin{enumerate}
    \item Perangkat lunak penyeimbang beban yang dievaluasi dibatasi pada NGINX \textit{Open Source} (versi stabil) dan HAProxy \textit{Community Edition}.
    \item Algoritma penyeimbang beban yang diuji meliputi empat varian: \textit{Round Robin} (RR), \textit{Least Connections} (LC), \textit{Weighted Round Robin} (WRR), dan \textit{Weighted Least Connections} (WLC).
    \item Lingkungan uji dibangun dengan kontainer Docker yang dikelola melalui Docker Compose pada sistem operasi Linux (kernel 6.x).
    \item Heterogenitas kapasitas kontainer \textit{backend} disimulasikan menggunakan fitur cgroups dengan perbandingan rasio kapasitas vCPU dan RAM sebesar 1:2:4 (masing-masing 0.5 vCPU/256 MB, 1.0 vCPU/512 MB, dan 2.0 vCPU/1024 MB) yang merepresentasikan tingkatan ukuran instansi standar (\textit{cloud compute tiers}: \textit{Small}, \textit{Medium}, \textit{Large}).
    \item Karakteristik pengujian beban dibagi menjadi dua kategori: \textit{endpoint} CPU-bound (komputasi intensif) dan \textit{endpoint} I/O-bound (penundaan I/O asinkron).
    \item Alat uji beban (\textit{load testing tool}) menggunakan Grafana k6 dengan metrik utama: \textit{throughput} (req/s), \textit{mean latency}, \textit{p95/p99 tail latency}, dan \textit{error rate}.
\end{enumerate}

\section{Tujuan Penelitian}
Tujuan yang hendak dicapai dalam pelaksanaan penelitian ini adalah:
\begin{enumerate}
    \item Menganalisis dan membandingkan performa NGINX dan HAProxy dalam menangani trafik HTTP pada lingkungan kontainer Docker berkapasitas heterogen.
    \item Mengevaluasi efektivitas algoritma penjadwalan berbobot (WRR dan WLC) dibandingkan algoritma tanpa bobot (RR dan LC) terhadap pencegahan lonjakan \textit{tail latency} dan \textit{CPU throttling}.
    \item Mengukur presisi distribusi trafik aktual untuk membuktikan kesesuaian antara alokasi bobot algoritma penyeimbang beban dengan kapasitas nyata kontainer \textit{backend}.
\end{enumerate}

\section{Manfaat Penelitian}
Hasil dari penelitian ini diharapkan dapat memberikan manfaat sebagai berikut:
\begin{enumerate}
    \item \textbf{Manfaat Teoretis}: Memberikan kontribusi ilmiah dan literatur empiris mengenai interaksi kausal antara mesin proksi lapis data-plane (arsitektur \textit{multi-process} NGINX vs \textit{event-driven} HAProxy) dengan mekanisme isolasi cgroups kernel Linux pada kluster kontainer heterogen.
    \item \textbf{Manfaat Praktis}: Menjadi panduan teknis bagi insinyur infrastruktur, praktisi \textit{DevOps}, dan arsitek sistem dalam memilih perangkat lunak penyeimbang beban serta menentukan algoritma distribusi yang optimal pada ekosistem komputasi awan berbasis kontainer.
\end{enumerate}

\section{Sistematika Penulisan}
Sistematika penulisan laporan proposal skripsi ini disusun dalam tiga bab utama sebagai berikut:
\begin{itemize}
    \item \textbf{BAB I PENDAHULUAN}: Memaparkan latar belakang masalah, perumusan masalah, batasan masalah, tujuan penelitian, manfaat penelitian, serta sistematika penulisan.
    \item \textbf{BAB II TINJAUAN PUSTAKA}: Menguraikan kajian literatur penelitian terdahulu yang relevan, dasar teori pendukung, serta kedudukan dan kebaruan penelitian.
    \item \textbf{BAB III METODOLOGI PENELITIAN}: Memaparkan alur kerangka penelitian, perancangan arsitektur topologi sistem, spesifikasi sumber daya heterogen, matriks skenario pengujian, instrumen evaluasi data, serta jadwal pelaksanaan penelitian.
\end{itemize}

% ==============================================================================
% BAB II: TINJAUAN PUSTAKA
% Format mengacu pada Pedoman Skripsi FST UNIBA Madura & Proposal Sempro Terakreditasi
% ==============================================================================

\chapter{TINJAUAN PUSTAKA}

\section{Penelitian Sebelumnya}
Kajian terhadap literatur terdahulu dilakukan guna memetakan perkembangan riset, mengevaluasi metodologi yang telah diterapkan, serta menemukan celah ilmiah (\textit{research gap}) yang menjadi dasar kebaruan penelitian ini. Penulis menelaah 11 karya ilmiah rujukan dalam rentang waktu lima tahun terakhir (2022--2026) yang mengkaji evaluasi kinerja penyeimbang beban (\textit{load balancing}), komparasi proksi NGINX dan HAProxy, implementasi algoritma penjadwalan statis maupun dinamis berbobot, serta dampaknya pada lingkungan fisik, mesin virtual, hingga kluster kontainer Docker. Pembahasan literatur terdahulu tersebut dikelompokkan ke dalam tiga domain utama sebagai berikut:

\subsection[Komparasi Mesin Penyeimbang Beban\\NGINX dan HAProxy]{Komparasi Mesin Penyeimbang Beban NGINX dan HAProxy}
Perbandingan langsung antara dua perangkat lunak penyeimbang beban sumber terbuka paling dominan di industri, yaitu NGINX dan HAProxy, telah menjadi fokus dari sejumlah penelitian empiris. \citet{nurzaman2022analisis} melakukan studi fondasional yang mengomparasikan performa \textit{head-to-head} antara NGINX dan HAProxy pada lingkungan kontainer Docker menggunakan algoritma statis \textit{Round Robin} (RR). Pengujian beban dilakukan menggunakan instrumen \texttt{httperf} pada rentang laju 200 hingga 1.000 permintaan per detik. Temuan empiris mereka mendemonstrasikan keandalan konkurensi HAProxy yang superior dengan tingkat kegagalan koneksi 0\% (tanpa \textit{connection timeout}), sedangkan NGINX mengalami kegagalan waktu tunggu rata-rata mencapai 739,8 permintaan pada laju trafik puncak. Namun, penelitian ini memiliki keterbatasan fundamental karena kontainer \textit{backend} dikonfigurasi secara seragam (homogen), hanya menguji satu algoritma tanpa bobot, serta tidak membedakan karakteristik beban komputasi prosesor versus operasi masukan/keluaran (I/O).

Karakteristik HAProxy dalam mengelola koneksi konkuren tinggi juga diinvestigasi oleh \citet{arini2025perbandingan} pada kluster basis data MySQL InnoDB menggunakan \textit{benchmark} Sysbench pada sistem operasi CentOS. Penelitian tersebut membuktikan bahwa efisiensi pengelolaan \textit{socket} dan manajemen \textit{event-loop} pada HAProxy menghasilkan latensi kueri yang sangat rendah dan utilisasi sumber daya yang stabil. Kendati demikian, fokus kajian mereka terbatas pada lapisan basis data (Layer 4 TCP), bukan pada penanganan trafik HTTP Lapisan Aplikasi (Layer 7). Sementara itu, \citet{mihai2023integrated} mengintegrasikan HAProxy dan NGINX pada sistem manajemen pembelajaran Moodle berbasis kontainer dengan beban kerja tinggi. Meskipun membuktikan ketahanan arsitektur gabungan dalam menopang lonjakan pengguna konkuren, kedua mesin proksi dipasang secara berseri (komplementer) di mana HAProxy bertindak sebagai gerbang terdepan dan NGINX sebagai peladen web perantara, bukan diuji secara bersaing (\textit{head-to-head}) pada kluster yang heterogen.

Pada domain evaluasi arsitektur NGINX, \citet{ikramsyah2025evaluating} menganalisis parameter \textit{Quality of Service} (QoS) meliputi \textit{throughput}, \textit{packet loss}, dan waktu respon pada penyeimbang beban NGINX. Hasil evaluasi menunjukkan bahwa stabilitas NGINX sangat dipengaruhi oleh volume trafik masuk dan pemilihan algoritma penjadwalan. Walaupun memberikan gambaran kinerja QoS pada NGINX, penelitian ini tidak menyertakan pembanding mesin HAProxy, menggunakan infrastruktur peladen yang homogen, dan tidak memisahkan pola beban kerja \textit{CPU-bound} dan \textit{I/O-bound}.

\subsection{Efektivitas Algoritma Penjadwalan Statis dan Berbobot}
Aspek krusial dalam rekayasa penyeimbang beban adalah penentuan algoritma distribusi trafik. Pada kondisi kluster dengan beban yang dinamis dan heterogen, algoritma statis seperti \textit{Round Robin} kerap memicu ketidakseimbangan antrean beban. \citet{zahir2025analisis} meneliti efisiensi peladen web NGINX dengan membandingkan algoritma \textit{Round Robin} dan \textit{Least Connections} (LC). Hasil eksperimen menunjukkan bahwa algoritma LC menghasilkan stabilitas waktu tanggap yang jauh lebih baik ketika jumlah koneksi klien meningkat, karena algoritma secara adaptif mengalokasikan trafik ke peladen dengan koneksi aktif paling sedikit. Namun, penelitian tersebut masih diuji pada mesin virtual homogen tanpa mempertimbangkan isolasi sumber daya kontainer maupun alokasi bobot kapasitas.

Keunggulan algoritma berbobot (\textit{weighted algorithms}) pada infrastruktur dengan disparitas kapasitas dibuktikan secara empiris oleh \citet{baeturohman2024analisis}. Mereka membandingkan kinerja algoritma \textit{Weighted Round Robin} (WRR) dan \textit{Weighted Least Connections} (WLC) pada NGINX menggunakan peladen virtual heterogen. Hasil penelitian mengungkapkan bahwa WLC memiliki ketahanan yang lebih tinggi dalam mempertahankan \textit{throughput} dan menekan latensi respon selama lonjakan trafik ekstrem dibandingkan WRR, karena WLC mengombinasikan faktor kapasitas teoretis (bobot) dengan jumlah koneksi riil yang sedang ditangani oleh masing-masing node. Penelitian serupa oleh \citet{rizqi2022analisis} pada infrastruktur \textit{Virtual Private Server} (VPS) juga menyimpulkan keunggulan WLC atas WRR dalam meminimalkan penumpukan antrean permintaan pada peladen berdaya rendah. Walaupun demikian, heterogenitas pada kedua penelitian tersebut dilakukan pada tingkat mesin virtual tradisional dan hanya membedakan kapasitas memori (RAM), sedangkan kuota siklus instruksi CPU tetap seragam dan tidak meneliti fenomena \textit{CPU throttling} pada kontainer Docker.

Penerapan algoritma berbobot pada lingkungan kontainer sempat dikaji oleh \citet{rahman2023perbandingan} yang menguji WRR dan WLC pada kluster Docker Swarm untuk layanan pengaliran (\textit{streaming}) video. Meskipun telah memanfaatkan teknologi kontainerisasi, fokus pengujian berorientasi pada metrik \textit{throughput} data video tanpa membedah interaksi mikroarsitektur \textit{data-plane} antara proses \textit{worker} NGINX dan \textit{thread} HAProxy terhadap pembatasan kuota \textit{Control Groups} (cgroups) Linux pada beban kerja mikroservis HTTP.

\subsection{Kinerja Penyeimbang Beban pada Infrastruktur Heterogen dan Kontainerisasi}
Tantangan nyata pada pusat data dan infrastruktur komputasi \textit{on-premise} modern adalah keberadaan peladen dengan spesifikasi perangkat keras yang heterogen. \citet{faridah2026performance} melakukan evaluasi komparatif enam algoritma \textit{load balancing} (\textit{Round Robin}, \textit{Least Connection}, \textit{Weighted Round Robin}, \textit{Weighted Least Connection}, \textit{IP Hash}, dan \textit{Random}) pada lingkungan peladen fisik \textit{on-premise} homogen dan heterogen menggunakan NGINX Plus dan Apache JMeter. Hasilnya membuktikan secara meyakinkan bahwa algoritma dinamis berbobot, khususnya WLC, menghasilkan \textit{throughput} tertinggi dan waktu tanggap terendah pada peladen heterogen, sekaligus mencegah terjadinya kelebihan beban (\textit{overload}) pada node berdaya rendah. Namun, penelitian ini diimplementasikan pada peladen fisik \textit{bare-metal}, belum menguji HAProxy, serta belum mengevaluasi mekanisme isolasi kontainer Docker dan metrik latensi persentil ekor (\textit{tail latency} p95/p99).

Urgensi penyeimbang beban pada lingkungan berdaya rendah juga dipertegas oleh \citet{ramadhani2026analysis} yang meneliti kinerja kluster peladen berbasis komputer papan tunggal (\textit{single-board computer}) Raspberry Pi untuk sistem informasi pariwisata. Penelitian ini menunjukkan bahwa perangkat keras berdaya rendah sangat rentan terhadap lonjakan konkurensi sehingga memerlukan distribusi beban yang adaptif. Sementara itu, \citet{sudiatmika2023round} meneliti algoritma \textit{Round Robin} pada pusat data nasional dan menekankan pentingnya evaluasi distribusi sebaran latensi respon guna menjaga mutu layanan publik. Namun, pengujian pada kedua studi tersebut masih terbatas pada peladen homogen dan tidak mengevaluasi algoritma berbobot secara komparatif antara NGINX dan HAProxy.

\subsection{Sintesis Posisi Penelitian}
Berdasarkan telaah mendalam terhadap 11 karya ilmiah rujukan di atas, dapat disimpulkan bahwa penelitian terdahulu masih terfragmentasi: studi yang membandingkan NGINX vs HAProxy pada Docker hanya menguji algoritma tanpa bobot pada node homogen, sedangkan studi yang menguji algoritma berbobot pada node heterogen dilakukan pada mesin virtual tradisional tanpa isolasi kuota prosesor cgroups dan tanpa pemisahan profil beban kerja. Oleh karena itu, penelitian ini memposisikan diri untuk mengisi kesenjangan riset tersebut secara komprehensif melalui evaluasi kinerja \textit{head-to-head} NGINX versus HAProxy dengan empat algoritma penjadwalan pada kluster Docker heterogen di bawah polaritas beban \textit{CPU-bound} dan \textit{I/O-bound}. Pemetaan komparatif posisi penelitian ini secara terperinci disajikan pada Bagian~\ref{sec:kedudukan_penelitian} (\tablename~\ref{tab:kedudukan_penelitian}).

\section{Dasar Teori}

\subsection{Konsep Penyeimbang Beban (Load Balancing)}
Penyeimbang beban (\textit{load balancer}) adalah komponen arsitektur sistem terdistribusi yang bertindak sebagai gerbang terdepan (\textit{reverse proxy}) dalam mengarahkan dan membagi lalu lintas permintaan HTTP/HTTPS masuk ke sekumpulan peladen komputasi \textit{backend} \citep{buke2026load, sakib2025dynamic, zahir2025analisis}. Beroperasi pada Lapisan Aplikasi (Layer 7 / HTTP) dari model OSI, \textit{load balancer} melakukan inspeksi terhadap paket data masuk, memutus koneksi TCP klien di sisi depan, dan menjalin koneksi baru ke peladen tujuan terpilih di sisi belakang.

Tujuan utama dari pengoperasian \textit{load balancer} meliputi pencegahan kelebihan beban (\textit{overload}) pada satu node tertentu, peningkatan ketersediaan layanan secara menyeluruh (\textit{high availability}), penyediaan mekanisme toleransi kesalahan (\textit{failover}), serta optimalisasi utilisasi sumber daya perangkat keras pada kluster peladen.

\subsection{Algoritma Penjadwalan Distribusi Beban}
Efektivitas distribusi trafik pada kluster peladen sangat dipengaruhi oleh algoritma penjadwalan yang diterapkan oleh penyeimbang beban \citep{buke2026load, sakib2025dynamic}:
\begin{enumerate}
    \item \textbf{Round Robin (RR)}: Algoritma statis paling mendasar yang mendistribusikan setiap permintaan baru secara berurutan melingkar (\textit{cyclic}) ke seluruh node peladen tanpa mempertimbangkan kapasitas sumber daya maupun jumlah beban aktif saat ini.
    \item \textbf{Least Connections (LC)}: Algoritma dinamis yang mengarahkan setiap permintaan baru ke peladen yang sedang melayani jumlah koneksi aktif paling sedikit pada saat permintaan diterima.
    \item \textbf{Weighted Round Robin (WRR)}: Pengembangan dari algoritma RR dengan menyematkan nilai bobot integer positif ($w_i$) pada masing-masing peladen ke-$i$. Peladen dengan nilai bobot lebih tinggi akan menerima proporsi permintaan yang lebih banyak secara deterministik. Probabilitas matematis penugasan permintaan ke peladen ke-$i$ dinyatakan dalam Persamaan (\ref{eq:wrr_bab2}):
    \begin{equation}
        P_i = \frac{w_i}{\sum_{j=1}^{n} w_j}
        \label{eq:wrr_bab2}
    \end{equation}
    di mana $P_i$ merupakan probabilitas teoretis penugasan beban ke peladen ke-$i$, dan $n$ menyatakan jumlah total peladen aktif di dalam kluster.
    \item \textbf{Weighted Least Connections (WLC)}: Algoritma dinamis tingkat lanjut yang memadukan nilai bobot kapasitas ($w_i$) dengan jumlah koneksi aktif riil ($C_i$) yang sedang ditangani oleh masing-masing peladen. Peladen terpilih ($S_{\text{selected}}$) adalah peladen yang memiliki rasio koneksi aktif terhadap bobot paling minimum, dirumuskan dalam Persamaan (\ref{eq:wlc_bab2}):
    \begin{equation}
        S_{\text{selected}} = \arg\min_{i} \left( \frac{C_i}{w_i} \right)
        \label{eq:wlc_bab2}
    \end{equation}
\end{enumerate}

\subsection{Arsitektur Data-Plane NGINX}
NGINX menerapkan model arsitektur berbasis banyak proses (\textit{multi-process}) asinkron yang digerakkan oleh peristiwa (\textit{event-driven}) \citep{faridah2026performance}. Struktur data-plane NGINX terdiri atas satu proses utama (\textit{master process}) dan beberapa proses pekerja (\textit{worker processes}) independen. Proses \textit{master} bertugas membaca konfigurasi dan mengelola siklus hidup proses \textit{worker}, sementara seluruh pemrosesan paket jaringan HTTP ditangani secara terisolasi oleh proses-proses \textit{worker}.

Setiap \textit{worker process} mengeksekusi satu \textit{event loop} terpisah yang memanfaatkan fasilitas pemanggilan sistem (\textit{system call}) Linux \texttt{epoll} non-pemblokir (\textit{non-blocking}) untuk menangani ribuan koneksi secara simultan. Namun, karena ruang alamat memori antar-proses di Linux bersifat terisolasi, sinkronisasi status koneksi aktif dan penentuan bobot algoritma antar-worker harus dijembatani melalui zona memori bersama (\textit{shared memory zones} / \texttt{shm}). Pada kondisi konkurensi ekstrem, mekanisme penguncian memori (\textit{mutex locking}) pada zona bersama tersebut dapat menimbulkan persaingan memori (\textit{memory contention}) yang berdampak pada penambahan latensi internal pemrosesan permintaan.

\subsection{Arsitektur Data-Plane HAProxy}
HAProxy didesain secara khusus sejak awal sebagai penyeimbang beban berkinerja sangat tinggi yang mengadopsi model arsitektur proses tunggal berbasis peristiwa (\textit{event-driven single process}) \citep{tarreau2024haproxy}. Berbeda dari NGINX, HAProxy meniadakan kebutuhan komunikasi antar-proses (IPC) dan penanganan \textit{mutex locking} memori bersama dalam mengelola status antrean koneksi.

Untuk penjadwalan dan penentuan peladen terbaik, HAProxy mengimplementasikan struktur data pohon biner elastis (\textit{Elastic Binary Tree} / \texttt{ebtree}). Struktur pohon biner berbasis kunci terindeks ini memungkinkan operasi penyisipan, pembaruan status beban koneksi, dan pencarian peladen berbeban terendah dieksekusi dengan kompleksitas waktu deterministik $O(1)$ untuk kasus terindeks dan $O(\log N)$ pada penataan antrean \citep{rawls2022load}. Karakteristik ini memberikan determinisme performa yang konsisten, khususnya saat mengelola antrean koneksi di tengah kondisi peladen mengalami degradasi sumber daya komputasi.

\subsection[Kontainerisasi Docker dan Linux cgroups]{Kontainerisasi Docker dan Linux Control Groups (cgroups)}
Docker mengoperasikan virtualisasi tingkat sistem operasi berbasis kontainer dengan mengandalkan dua fitur inti kernel Linux, yaitu \textit{Namespaces} untuk isolasi ruang pandang sistem dan \textit{Control Groups} (cgroups v2) untuk pembatasan alokasi sumber daya fisik \citep{khan2026decomposing, gurramkonda2026docker}.

Pada subsistem manajemen prosesor, kernel Linux menggunakan penjadwal \textit{Completely Fair Scheduler} (CFS) \citep{shetty2025cpu}. Pembatasan kuota CPU pada kontainer Docker diatur melalui dua parameter cgroups: periode kuota (\texttt{cpu.cfs\_}\allowbreak\texttt{period\_us}, berstandar $100.000\,\mu\text{s}$ atau $100\,\text{ms}$) dan batas alokasi kuota eksekusi (\texttt{cpu.cfs\_}\allowbreak\texttt{quota\_us}). Ketika proses di dalam kontainer mengonsumsi waktu prosesor yang melebihi batas kuota dalam satu siklus, kernel Linux secara otomatis memberlakukan pemotongan jatah waktu CPU (\textit{CPU throttling}) \citep{shetty2025cpu}. Pemotongan ini membekukan eksekusi proses kontainer hingga siklus cgroups berikutnya, yang memicu lonjakan antrean soket jaringan dan latensi secara tajam pada sisi klien.

\subsection[Karakteristik Beban Kerja CPU-Bound dan I/O-Bound]{Karakteristik Beban Kerja Komputasi (CPU-Bound vs. I/O-Bound)}
Layanan web modern melayani spektrum beban kerja dengan karakteristik komputasi yang sangat kontras:
\begin{enumerate}
    \item \textbf{Beban Intensif Prosesor (CPU-Bound)}: Beban kerja yang didominasi oleh operasi aritmatika, kalkulasi algoritma, atau pemrosesan data (seperti komputasi deret rekursif Fibonacci). Karakteristik beban ini memforsir siklus instruksi CPU secara kontinu, sehingga sangat cepat menguras kuota waktu CFS dan memicu \textit{CPU throttling} pada kontainer berkapasitas rendah.
    \item \textbf{Beban Intensif Masukan/Keluaran (I/O-Bound)}: Beban kerja yang didominasi oleh operasi tunggu jaringan atau akses penyimpanan (seperti kueri basis data atau jeda simulasi I/O non-pemblokir). Karakteristik beban ini tidak menguras siklus CPU, melainkan menguji kapasitas penanganan antrean soket jaringan konkuren dan ketahanan memori penyeimbang beban saat ribuan koneksi tertahan dalam kondisi terbuka (\textit{keep-alive}).
\end{enumerate}
Pemisahan kedua karakteristik ini mutlak diperlukan dalam evaluasi komparatif guna memetakan korelasi kausal antara logika algoritma penjadwalan dan sifat konsumsi sumber daya sistem.

\subsection{Pengujian Beban Terdistribusi dengan Grafana k6}
Grafana k6 adalah instrumen pengujian beban (\textit{load testing}) modern bersumber terbuka yang dibangun menggunakan bahasa pemrograman Go \citep{k6docs2024, sudiatmika2023round}. Berbeda dengan instrumen berbasis JVM seperti Apache JMeter yang membebankan \textit{thread} sistem operasi untuk setiap pengguna virtual, k6 mengeksekusi model konkurensi \textit{goroutine} ringan yang sangat efisien dalam penggunaan memori dan siklus prosesor mesin penguji.

Grafana k6 mendukung otomasi pengujian berbasis skrip JavaScript ES6, tahapan pembebanan dinamis (\textit{ramping stages}), serta pencatatan metrik performa beresolusi tinggi hingga skala mikrodetik untuk persentil latensi ekor.

\subsection{Metrik Evaluasi Kinerja Sistem}
Evaluasi kinerja komparatif sistem diukur menggunakan enam parameter kuantitatif terstandarisasi:
\begin{enumerate}
    \item \textbf{Throughput}: Kuantitas permintaan HTTP yang berhasil diproses dan diselesaikan secara tuntas per satuan waktu, dinyatakan dalam \textit{Requests Per Second} (RPS).
    \item \textbf{Latensi Rerata (Mean Latency)}: Rata-rata waktu tempuh bolak-balik (\textit{Round Trip Time} / RTT) yang dihitung sejak paket permintaan HTTP dikirimkan oleh mesin uji hingga paket respon pertama diterima kembali secara utuh.
    \item \textbf{Tail Latency (Persentil p95 dan p99)}: Waktu respon pada titik persentil ke-95 dan ke-99 dari total distribusi permintaan yang berhasil diselesaikan \citep{wodislawski2026ponytail, li2025tailbench}. Metrik ini krusial dalam mendeteksi degradasi performa ekstrem yang kerap tersembunyi oleh metrik rata-rata pada sistem heterogen. Rasio amplifikasi latensi ekor dinyatakan dalam Persamaan (\ref{eq:rtail_bab2}):
    \begin{equation}
        R_{\text{tail}} = \frac{L_{p99}}{L_{\text{mean}}}
        \label{eq:rtail_bab2}
    \end{equation}
    di mana $L_{p99}$ adalah waktu respon persentil ke-99 dan $L_{\text{mean}}$ adalah waktu respon rata-rata. Nilai $R_{\text{tail}} > 1$ mengindikasikan tingkat ketimpangan latensi ekstrem yang dialami sebagian pengguna.
    \item \textbf{Tingkat Kegagalan (Error Rate)}: Persentase kegagalan pemrosesan permintaan dibandingkan total permintaan yang dikirimkan, mencakup respon kesalahan peladen (kode status HTTP 5xx) dan kegagalan waktu tunggu soket (\textit{connection timeout}):
    \begin{equation}
        E = \frac{N_{\text{gagal}}}{N_{\text{total}}} \times 100\%
        \label{eq:error_rate_bab2}
    \end{equation}
    \item \textbf{Audit Presisi Distribusi Beban (MAPE)}: Derajat kedekatan antara persentase jumlah permintaan aktual yang diterima oleh masing-masing kontainer (\textit{log audit}) terhadap rasio kapasitas kuota teoretisnya, diukur menggunakan formulasi \textit{Mean Absolute Percentage Error} (MAPE):
    \begin{equation}
        \text{MAPE} = \frac{1}{n} \sum_{i=1}^{n} \left| \frac{R_i^{\text{aktual}} - R_i^{\text{teoretis}}}{R_i^{\text{teoretis}}} \right| \times 100\%
        \label{eq:mape_bab2}
    \end{equation}
    di mana $R_i^{\text{aktual}}$ adalah persentase trafik riil yang tercatat pada log akses kontainer ke-$i$, dan $R_i^{\text{teoretis}}$ adalah persentase bobot alokasi idealnya.
\end{enumerate}

\section{Kedudukan Penelitian}
\label{sec:kedudukan_penelitian}

\subsection{Sintesis dan Matriks Pemetaan Posisi Penelitian}
Berdasarkan telaah mendalam terhadap 11 penelitian rujukan utama dalam rentang tahun 2022 hingga 2026, belum ditemukan studi yang secara komprehensif mengintegrasikan perbandingan mesin penyeimbang beban, matriks algoritma berbobot, kluster kontainer Docker berkapasitas heterogen, pemisahan polaritas beban kerja, serta pengukuran metrik latensi persentil. Matriks pemetaan posisi penelitian ini dibandingkan dengan literatur rujukan disajikan pada \tablename\ \ref{tab:kedudukan_penelitian}.

\begin{table}[H]
    \centering
    \caption{Matriks Pemetaan Posisi dan Kebaruan Penelitian}
    \label{tab:kedudukan_penelitian}
    \small
    \begin{tabularx}{\textwidth}{|>{\raggedright\arraybackslash}p{2.8cm}|*{5}{>{\centering\arraybackslash}X|}}
        \hline
        \rowcolor[gray]{0.85}
        \textbf{Peneliti \& Tahun} & \textbf{NGINX vs HA} & \textbf{4 Algoritma} & \textbf{Docker Het.} & \textbf{Beban CPU/IO} & \textbf{Tail Latency} \\
        \hline
        \citet{faridah2026performance} & Tidak & Sebagian & Tidak & Tidak & Tidak \\
        \hline
        \citet{ramadhani2026analysis}  & Tidak & Tidak    & Tidak & Tidak & Tidak \\
        \hline
        \citet{zahir2025analisis}      & Tidak & Tidak    & Tidak & Tidak & Tidak \\
        \hline
        \citet{ikramsyah2025evaluating}& Tidak & Tidak    & Tidak & Tidak & Tidak \\
        \hline
        \citet{arini2025perbandingan}  & Tidak & Tidak    & Tidak & Khusus DB & Tidak \\
        \hline
        \citet{baeturohman2024analisis}& Tidak & Sebagian & Tidak (VM) & Tidak & Tidak \\
        \hline
        \citet{rahman2023perbandingan} & Tidak & Sebagian & Tidak & Streaming & Tidak \\
        \hline
        \citet{mihai2023integrated}    & Berurutan & Tidak & Sebagian & Moodle & Tidak \\
        \hline
        \citet{sudiatmika2023round}    & Tidak & Tidak    & Tidak & Tidak & Sebagian \\
        \hline
        \citet{rizqi2022analisis}      & Tidak & Sebagian & Sebagian (VPS) & Tidak & Tidak \\
        \hline
        \citet{nurzaman2022analisis}   & Ya    & Tidak (RR) & Tidak & Tidak & Tidak \\
        \hline
        \rowcolor[gray]{0.92}
        \textbf{Penelitian Ini}        & \textbf{Ya} & \textbf{Ya (4 Alg)} & \textbf{Ya (cgroups)} & \textbf{Ya (Terpisah)} & \textbf{Ya (p95 \& p99)} \\
        \hline
    \end{tabularx}

\end{table}

\subsection{Identifikasi Kesenjangan Riset (Research Gaps)}
Berdasarkan telaah komparatif pada \tablename\ \ref{tab:kedudukan_penelitian}, teridentifikasi empat kesenjangan riset (\textit{research gaps}) fundamental yang menjadi landasan orisinalitas penelitian ini:
\begin{enumerate}
    \item \textbf{Asumsi Homogenitas pada Docker vs. Heterogenitas Parsial pada Mesin Virtual}: Penelitian terdahulu yang membandingkan performa NGINX dan HAProxy secara langsung di lingkungan Docker \citep{nurzaman2022analisis} seluruhnya mengonfigurasi kontainer berkapasitas identik (homogen) dan hanya menguji algoritma tanpa bobot (\textit{Round Robin}). Sebaliknya, penelitian yang menguji peladen heterogen dengan algoritma berbobot \citep{baeturohman2024analisis, rizqi2022analisis} dilakukan di atas mesin virtual tradisional (VirtualBox/VPS) tanpa memperhitungkan mekanisme isolasi kuota prosesor berbasis \textit{Control Groups} (cgroups) kernel Linux. Fenomena kritis di mana kontainer kecil mengalami pemotongan waktu eksekusi (\textit{CPU throttling}) saat jatah kuota habis belum pernah dipetakan secara komparatif antara NGINX dan HAProxy.
    \item \textbf{Kebutaan terhadap Dikotomi Karakteristik Beban Kerja (CPU-Bound vs. I/O-Bound)}: Sebagian besar literatur membebani kluster peladen web menggunakan profil trafik tunggal yang seragam. Pendekatan ini mengaburkan korelasi kausal antara logika pemilihan peladen dan sifat konsumsi sumber daya. Beban kerja \textit{CPU-bound} menguras siklus instruksi prosesor dan memicu penalti penjadwalan CFS kernel, sedangkan beban \textit{I/O-bound} menahan soket jaringan terbuka tanpa konsumsi CPU tinggi. Belum ada penelitian yang membedah ketahanan mikroarsitektur antrean kedua penyeimbang beban terhadap polaritas karakteristik beban ini pada kluster heterogen.
    \item \textbf{Fragmentasi Matriks Pengujian dan Absennya Audit Presisi Proporsi Aktual}: Literatur yang menguji komparasi NGINX vs HAProxy belum menerapkan algoritma berbobot secara lengkap (WRR dan WLC), sedangkan literatur yang mengkaji algoritma berbobot hanya meneliti satu jenis perangkat lunak (NGINX). Selain itu, belum ada penelitian empiris yang melakukan audit kuantitatif terhadap \textit{log} akses kontainer guna memvalidasi apakah proporsi distribusi trafik aktual benar-benar presisi mengikuti rasio alokasi cgroups saat terjadi kejenuhan beban.
    \item \textbf{Bias Penggunaan Nilai Rata-rata (\textit{Average Bias})}: Mayoritas riset terdahulu hanya menyajikan nilai rata-rata (\textit{mean latency}). Pada kluster dengan disparitas kapasitas, nilai rata-rata memberikan ilusi performa karena waktu respon cepat dari kontainer berdaya tinggi menutupi kelambatan ekstrem yang dialami permintaan yang terjebak di antrean kontainer kecil. Oleh sebab itu, evaluasi berbasis persentil ke-95 (p95) dan ke-99 (p99) menjadi mutlak untuk menangkap degradasi nyata yang dialami pengguna.
\end{enumerate}

\subsection{Pernyataan Kebaruan (Novelty Statement)}
Berpijak pada empat kesenjangan riset tersebut, kebaruan (\textit{novelty}) yang dihadirkan dalam penelitian ini dirumuskan sebagai berikut:
\begin{quote}
    \textit{"Evaluasi komparatif komprehensif performa lapisan data-plane NGINX versus HAProxy dengan matriks empat algoritma penjadwalan (Round Robin, Least Connections, Weighted Round Robin, dan Weighted Least Connections) pada kluster kontainer Docker dengan kapasitas asimetris berbasis cgroups Linux, melalui pemisahan tegas karakteristik beban kerja CPU-bound dan I/O-bound, pengukuran metrik latensi ekor (tail latency p95 dan p99) menggunakan Grafana k6, serta audit presisi proporsi distribusi beban aktual (MAPE)."}
\end{quote}

\subsection{Kerangka Berpikir Penelitian}
Kerangka berpikir penelitian disusun secara sistematis untuk menghubungkan latar belakang permasalahan, kesenjangan literatur, perancangan matriks eksperimen, hingga luaran yang ditargetkan. Alur kerangka berpikir penelitian ini diilustrasikan pada \figurename\ \ref{fig:kerangka_berpikir}.

\begin{figure}[H]
    \centering
    \singlespacing
    \begin{tikzpicture}[
        node distance=0.35cm and 1.0cm,
        box/.style={
            rectangle,
            draw=black!80,
            fill=gray!10,
            thick,
            rounded corners=3pt,
            text width=12.5cm,
            align=left,
            inner sep=3.5pt,
            font=\footnotesize
        },
        arrow/.style={
            -{Stealth[length=2.0mm]},
            thick,
            draw=black!80
        }
    ]
        % Node 1: Masalah Utama
        \node[box] (masalah) {
            \textbf{1. Masalah Utama dan Tantangan Nyata:}\\
            Asimetri sumber daya (\textit{hardware heterogeneity}) pada lingkungan kontainer Docker memicu \textit{CPU throttling} pada kontainer kecil saat menerima beban tinggi, menimbulkan lonjakan latensi ekor (\textit{tail latency}) dan kegagalan koneksi jika beban tidak didistribusikan secara proporsional.
        };

        % Node 2: Kesenjangan Riset
        \node[box, below=of masalah] (gap) {
            \textbf{2. Identifikasi Kesenjangan Riset (Research Gaps):}\\
            $\bullet$ Evaluasi NGINX vs HAProxy di Docker terpaku pada kluster homogen tanpa algoritma bobot.\\
            $\bullet$ Belum ada pengujian terpisah antara karakteristik beban \textit{CPU-bound} versus \textit{I/O-bound}.\\
            $\bullet$ Minimnya evaluasi metrik \textit{tail latency} (p95 dan p99) serta absennya audit presisi distribusi beban.
        };

        % Node 3: Solusi Eksperimental
        \node[box, below=of gap] (solusi) {
            \textbf{3. Rancangan Eksperimental Komparatif:}\\
            $\bullet$ \textbf{Matriks Komparasi}: 2 Mesin Proksi (NGINX vs HAProxy) $\times$ 4 Algoritma (RR, LC, WRR, WLC).\\
            $\bullet$ \textbf{Lingkungan Backend}: Kluster Docker Heterogen 1:2:4 core CPU dan RAM via cgroups Linux.\\
            $\bullet$ \textbf{Pemisahan Beban}: Beban \textit{CPU-bound} (Fibonacci rekursif) dan \textit{I/O-bound} (jeda non-pemblokir).
        };

        % Node 4: Pengujian Beban
        \node[box, below=of solusi] (uji) {
            \textbf{4. Pengujian Beban dan Pengukuran Metrik (Grafana k6):}\\
            $\bullet$ Simulasi beban bertingkat (\textit{ramping}) dari 50 hingga 1.000 pengguna virtual (\textit{Virtual Users}).\\
            $\bullet$ Ekstraksi metrik kuantitatif: \textit{Throughput} (RPS), Rerata Latensi, \textit{Tail Latency} (p95 dan p99).\\
            $\bullet$ Pengukuran \textit{Error Rate} (HTTP 5xx \& \textit{timeout}) dan Audit Presisi Distribusi Beban (MAPE).
        };

        % Node 5: Luaran Penelitian
        \node[box, below=of uji] (output) {
            \textbf{5. Luaran yang Diharapkan (\textit{Expected Outcomes}):}\\
            $\bullet$ Pemetaan kausal empiris interaksi mikroarsitektur antrean NGINX vs HAProxy terhadap \textit{cgroups}.\\
            $\bullet$ Panduan rekomendasi konfigurasi algoritma dan penyeimbang beban optimal untuk arsitektur heterogen.\\
            $\bullet$ Data patokan (\textit{benchmark}) kinerja latensi ekor web service pada lingkungan kluster kontainer.
        };

        % Arrows
        \draw[arrow] (masalah) -- (gap);
        \draw[arrow] (gap) -- (solusi);
        \draw[arrow] (solusi) -- (uji);
        \draw[arrow] (uji) -- (output);
    \end{tikzpicture}
    \caption{Kerangka Berpikir Penelitian}
    \label{fig:kerangka_berpikir}
\end{figure}

Berdasarkan \figurename\ \ref{fig:kerangka_berpikir}, alur penalaran penelitian diawali dari permasalahan asimetri kapasitas pada kluster kontainer Docker yang berpotensi memicu degradasi latensi ekstrem saat beban kerja tinggi. Berangkat dari identifikasi empat kesenjangan riset pada literatur ilmiah terkini, dirancanglah sebuah kerangka eksperimen komparatif terpadu yang memadukan dua mesin proksi utama (NGINX dan HAProxy) dengan empat algoritma distribusi beban pada kluster \textit{backend} heterogen dengan rasio 1:2:4 berbasis cgroups Linux. Eksperimen diuji di bawah dua polaritas beban kerja (CPU-bound dan I/O-bound) menggunakan instrumen Grafana k6 untuk mengekstrak data performa secara presisi, sehingga menghasilkan luaran berupa bukti empiris kausal, rekomendasi konfigurasi optimal, dan basis data patokan kinerja sistem.

% ==============================================================================
% BAB III: METODOLOGI PENELITIAN
% ==============================================================================

\chapter{METODOLOGI PENELITIAN}

\section{Kerangka Penelitian}
Penelitian ini dilaksanakan secara bertahap dan sistematis mengikuti metodologi eksperimen kuantitatif. Alur tahapan penelitian dirancang secara terstruktur untuk menjawab setiap rumusan masalah, mulai dari studi kepustakaan, perancangan arsitektur, konfigurasi lingkungan heterogen, pelaksanaan pengujian stres, pengumpulan dan audit data, hingga penarikan kesimpulan. Kerangka metodologi penelitian ini disajikan dalam bentuk diagram alir pada \figurename\ \ref{fig:kerangka_penelitian}.

\begin{figure}[H]
    \centering
    \begin{tikzpicture}[
        node distance=0.6cm and 1.0cm,
        startstop/.style={rectangle, rounded corners=5pt, minimum width=3.6cm, minimum height=0.55cm, text centered, draw=black, fill=gray!20, font=\footnotesize\bfseries, inner sep=3pt},
        process/.style={rectangle, minimum width=8.6cm, minimum height=0.65cm, text centered, text width=8.2cm, draw=black, fill=white, font=\footnotesize, inner sep=3pt},
        decision/.style={diamond, aspect=2.4, text centered, text width=3.2cm, draw=black, fill=gray!10, font=\footnotesize, inner sep=2pt},
        arrow/.style={-{Stealth[scale=0.9]}, thick}
    ]
        % Nodes
        \node (start) [startstop] {Mulai};
        
        \node (lit) [process, below=0.36cm of start] {%
            \textbf{1. Studi Literatur \& Analisis Kesenjangan Riset}\\
            Kajian teori penyeimbang beban, karakteristik cgroups kernel Linux, dan sintesis 11 riset komparatif terdahulu
        };
        
        \node (design) [process, below=0.36cm of lit] {%
            \textbf{2. Perancangan Sistem \& Lingkungan Eksperimen}\\
            Perancangan topologi Docker Bridge, penentuan rasio kapasitas 1:2:4, dan implementasi endpoint CPU-bound \& I/O-bound
        };
        
        \node (impl) [process, below=0.36cm of design] {%
            \textbf{3. Konfigurasi Penyeimbang Beban \& Telemetri}\\
            Penyusunan profil NGINX dan HAProxy (RR, LC, WRR, WLC) serta integrasi instrumen metrik k6, cAdvisor, dan Prometheus
        };
        
        \node (test) [process, below=0.36cm of impl] {%
            \textbf{4. Pelaksanaan Uji Beban Bertahap (Stress Testing)}\\
            Eksekusi 16 matriks skenario pengujian dengan variasi beban konkurensi bertahap (100 s.d. 1.000 VUs)
        };
        
        \node (valid) [decision, below=0.36cm of test] {%
            Data Metrik \& Log Valid/Lengkap?
        };
        
        \node (analisis) [process, below=0.38cm of valid] {%
            \textbf{5. Analisis Komparatif \& Audit Distribusi Beban}\\
            Evaluasi throughput, latensi rata-rata, tail latency (p95 \& p99), fenomena CPU throttling, dan audit presisi beban aktual
        };
        
        \node (selesai) [startstop, below=0.36cm of analisis] {Kesimpulan \& Selesai};

        % Arrows
        \draw [arrow] (start) -- (lit);
        \draw [arrow] (lit) -- (design);
        \draw [arrow] (design) -- (impl);
        \draw [arrow] (impl) -- (test);
        \draw [arrow] (test) -- (valid);
        \draw [arrow] (valid) -- node[anchor=east, font=\scriptsize\bfseries] {Ya} (analisis);
        \draw [arrow] (analisis) -- (selesai);
        
        % Loop back if data invalid
        \draw [arrow] (valid.east) -- ++(1.0,0) node[anchor=south, font=\scriptsize\bfseries] {Tidak (Ulangi)} |- (test.east);
    \end{tikzpicture}
    \vspace{-4pt}
    \caption{Diagram Alir Kerangka Metodologi Penelitian}
    \label{fig:kerangka_penelitian}
\end{figure}

\clearpage

Secara terperinci, tahapan-tahapan yang ditempuh dalam penelitian ini dirancang sebagai berikut:
\begin{enumerate}
    \item \textbf{Tahap Studi Literatur dan Analisis Kesenjangan Riset}: Mengkaji landasan teoretis arsitektur \textit{reverse proxy} pada lapisan aplikasi, mekanisme penjadwalan \textit{Completely Fair Scheduler} (CFS) pada isolasi kuota Linux \textit{Control Groups} (cgroups v2), serta menganalisis kesenjangan riset dari publikasi rujukan terdahulu guna memvalidasi orisinalitas perumusan masalah asimetri sumber daya.
    \item \textbf{Tahap Perancangan Sistem dan Topologi Uji}: Merancang lingkungan uji komputasi \textit{bare-metal} pada laptop MSI Thin GF63 12UCX (Ubuntu 24.04 LTS) dengan partisi isolasi CPU (\textit{CPU pinning}) berbasis arsitektur hibrida P-Cores dan E-Cores, merancang jaringan privat \textit{Docker Bridge Network} (\texttt{skripsi-net}), merumuskan alokasi kuota komputasi heterogen berasio kapasitas 1:2:4 (0,5 vCPU/256 MB, 1,0 vCPU/512 MB, dan 2,0 vCPU/1024 MB), serta mengonstruksi beban kerja sintetis kontras (\textit{CPU-bound} dan \textit{I/O-bound}).
    \item \textbf{Tahap Konfigurasi dan Instrumentasi Telemetri}: Menyusun profil konfigurasi penyeimbang beban NGINX dan HAProxy untuk empat algoritma penjadwalan (RR, LC, WRR, dan WLC), mengekualisasikan batas konkurensi serta parameter kernel (\textit{sysctl}), memformat rekaman log akses presisi tinggi, serta mengintegrasikan instrumen telemetri cAdvisor dan Prometheus untuk memantau pemotongan waktu prosesor (\textit{CPU throttling}).
    \item \textbf{Tahap Pelaksanaan Pengujian Beban Stres}: Mengeksekusi 16 skenario pengujian faktorial penuh menggunakan perkakas Grafana k6 dengan profil konkurensi bertahap (\textit{ramping Virtual Users}) dari 100 hingga 1.000 pengguna virtual selama durasi sepuluh menit per iterasi, di mana masing-masing skenario diulang tiga kali replikasi independen guna menjamin reliabilitas empiris.
    \item \textbf{Tahap Ekstraksi Metrik, Audit Log, dan Analisis Komparatif}: Menghimpun data metrik kinerja sisi klien (\textit{throughput}, rerata latensi, latensi ekor p95 dan p99) dan telemetri utilisasi sumber daya sisi peladen, dilanjutkan dengan audit presisi proporsi distribusi trafik aktual terhadap target ideal melalui kalkulasi MAPE serta uji signifikansi komparatif statistik.
    \item \textbf{Tahap Sintesis Temuan dan Penarikan Kesimpulan}: Mengkaji korelasi kausal antara interaksi antrean penyeimbang beban terhadap dinamika penjadwalan kernel Linux, merumuskan matriks rekomendasi konfigurasi optimal bagi arsitektur heterogen, serta menyusun patokan empiris sebagai pedoman perancangan sistem berbasis kontainer.
\end{enumerate}

\clearpage

\section{Perancangan Arsitektur dan Topologi Sistem}
Arsitektur sistem dibangun di atas lingkungan kontainerisasi Docker yang diisolasi dalam satu jaringan privat (\textit{Docker Bridge Network}). Penyeimbang beban (NGINX atau HAProxy) diposisikan di baris depan (\textit{reverse proxy}) untuk menerima permintaan trafik HTTP dari mesin penguji (\textit{load generator}), kemudian meneruskannya ke tiga kontainer peladen aplikasi \textit{backend}, seperti diilustrasikan pada \figurename\ \ref{fig:topologi_sistem}.

\begin{figure}[H]
    \centering
    \singlespacing
    \begin{tikzpicture}[
        scale=0.78, transform shape,
        font=\scriptsize,
        >=Stealth,
        node distance=0.20cm and 0.8cm,
        box/.style={
            draw=black!80,
            rounded corners=3pt,
            fill=white,
            thick,
            inner sep=2.5pt,
            align=center
        },
        clientbox/.style={
            box,
            fill=gray!10,
            text width=3.3cm
        },
        lbbox/.style={
            box,
            fill=gray!15,
            text width=2.5cm
        },
        bebox/.style={
            box,
            fill=white,
            text width=4.0cm
        }
    ]
        % -------------------------------------------------------------
        % 1. CLIENT DOMAIN (Sisi Kiri) - Diisolasi pada E-Cores (CPU 8-11)
        % -------------------------------------------------------------
        \node[clientbox] (client) {
            \textbf{Domain Pembangkit Beban}\\[1pt]
            \textit{Isolasi CPU: E-Cores (CPU 8--11)}\\
            (\texttt{taskset -c 8-11})\\[2pt]
            \rule{2.8cm}{0.3pt}\\[2pt]
            \textbf{Grafana k6}\\
            \textit{Load Generator}\\
            (100 s.d. 1.000 VUs)
        };

        % -------------------------------------------------------------
        % 2. SERVER NODE - REVERSE PROXY (Tengah)
        % -------------------------------------------------------------
        \node[lbbox, right=1.6cm of client] (lb) {
            \textbf{Reverse Proxy}\\[1pt]
            \textbf{NGINX / HAProxy}\\
            Port: \texttt{80}\\[2pt]
            \rule{2.2cm}{0.3pt}\\[2pt]
            \textit{Algoritma:}\\
            RR, LC, WRR, WLC
        };

        % -------------------------------------------------------------
        % 3. BACKEND CONTAINERS (Sisi Kanan)
        % -------------------------------------------------------------
        \node[bebox, right=1.1cm of lb] (be2) {
            \textbf{\texttt{web-backend-2} (Tier Medium)}\\
            1.0 vCPU $\mid$ 512 MB RAM $\mid$ Port 8080\\
            Rasio Kapasitas: 2 $\mid$ Bobot: $w_2 = 2$
        };
        \node[bebox, above=0.16cm of be2] (be1) {
            \textbf{\texttt{web-backend-1} (Tier Small)}\\
            0.5 vCPU $\mid$ 256 MB RAM $\mid$ Port 8080\\
            Rasio Kapasitas: 1 $\mid$ Bobot: $w_1 = 1$
        };
        \node[bebox, below=0.16cm of be2] (be3) {
            \textbf{\texttt{web-backend-3} (Tier Large)}\\
            2.0 vCPU $\mid$ 1024 MB RAM $\mid$ Port 8080\\
            Rasio Kapasitas: 4 $\mid$ Bobot: $w_3 = 4$
        };

        % -------------------------------------------------------------
        % 4. BOUNDING BOXES (FRAMES LINGKUNGAN)
        % -------------------------------------------------------------
        % Docker Network Frame
        \node[draw=black!70, dashed, rounded corners=3pt, inner xsep=5pt, inner ysep=5pt,
              fit=(be1) (be3)] (dockernet) {};
        \node[anchor=south west, font=\tiny\bfseries, inner sep=1pt] at (dockernet.north west) {Docker Network (\texttt{skripsi-net}: 172.28.0.0/16)};

        % Server Node Frame (P-Cores)
        \node[draw=black!90, thick, rounded corners=4pt, inner xsep=6pt, inner ysep=7pt,
              fit=(lb) (dockernet)] (servernode) {};
        \node[anchor=south east, font=\scriptsize\bfseries, inner sep=1pt] at (servernode.north east) {Domain Peladen: P-Cores (CPU 0--7, \texttt{cpuset})};

        % Outer Frame: Host Fisik (MSI Thin GF63 12UCX)
        \node[draw=black!60, very thick, rounded corners=6pt, inner xsep=8pt, inner ysep=12pt,
              fit=(client) (servernode)] (hostframe) {};
        \node[anchor=north west, font=\tiny\bfseries, inner sep=2pt] at (hostframe.north west) {Host Fisik: Laptop MSI Thin GF63 12UCX (Ubuntu 24.04 LTS, RAM 16 GB, SSD 100 GB)};

        % -------------------------------------------------------------
        % 5. ARROWS & CONNECTIONS
        % -------------------------------------------------------------
        % Local Loopback Connection
        \draw[->, thick] (client.east) -- node[above, font=\tiny\bfseries, align=center] {Trafik HTTP\\(Port 80)} 
                                          node[below, font=\tiny, align=center] {Local Loopback\\(\texttt{127.0.0.1})} (lb.west);

        % Distribution to Backends
        \draw[->, thick] (lb.east) -- ++(0.5,0) |- (be1.west);
        \draw[->, thick] (lb.east) -- (be2.west);
        \draw[->, thick] (lb.east) -- ++(0.5,0) |- (be3.west);

    \end{tikzpicture}
    \caption{Arsitektur dan Topologi Pengujian Mandiri Kluster Kontainer Heterogen}
    \label{fig:topologi_sistem}
\end{figure}

Topologi kluster kontainer \textit{backend} dirancang dengan konfigurasi kapasitas sumber daya yang heterogen menggunakan fitur pembatasan kuota \textit{Control Groups} (cgroups v2) pada Docker Compose, sebagaimana ditunjukkan pada \tablename\ \ref{tab:spesifikasi_backend}.

\begin{table}[H]
    \centering
    \caption{Spesifikasi Alokasi Sumber Daya Kontainer Backend Heterogen}
    \label{tab:spesifikasi_backend}
    \small
    \begin{tabularx}{\textwidth}{|p{3.8cm}|*{4}{>{\centering\arraybackslash}X|}}
        \hline
        \rowcolor[gray]{0.85}
        \textbf{Nama Kontainer} & \textbf{Alokasi vCPU} & \textbf{Batas RAM} & \textbf{Rasio Kapasitas} & \textbf{Bobot Ideal ($w_i$)} \\
        \hline
        \texttt{web-backend-1} & 0.5 vCPU & 256 MB  & 1 & 1 \\
        \hline
        \texttt{web-backend-2} & 1.0 vCPU & 512 MB  & 2 & 2 \\
        \hline
        \texttt{web-backend-3} & 2.0 vCPU & 1024 MB & 4 & 4 \\
        \hline
    \end{tabularx}
\end{table}

Guna menjamin keabsahan data empiris dan mencegah bias pengukuran akibat perebutan sumber daya prosesor antara pembangkit beban (\textit{load generator}) dan peladen yang diuji (\textit{server under test}), lingkungan eksperimen diimplementasikan pada perangkat keras \textit{bare-metal} laptop MSI Thin GF63 12UCX dengan partisi sistem operasi mandiri Ubuntu 24.04 LTS. Eksperimen memanfaatkan arsitektur prosesor hibrida Intel Core i5-12450H melalui mekanisme pemisahan core secara terisolasi (\textit{CPU affinity / CPU pinning}), sebagaimana dirinci pada \tablename\ \ref{tab:spesifikasi_lingkungan_uji}.

\begin{table}[H]
    \centering
    \caption{Spesifikasi Perangkat Keras dan Lingkungan Pengujian Mandiri}
    \label{tab:spesifikasi_lingkungan_uji}
    \small
    \begin{tabularx}{\textwidth}{|p{3.2cm}|p{4.6cm}|>{\raggedright\arraybackslash}X|}
        \hline
        \rowcolor[gray]{0.85}
        \textbf{Komponen Sistem} & \textbf{Spesifikasi Rinci} & \textbf{Alokasi dan Peran Pengujian} \\
        \hline
        \textbf{Model Perangkat} & MSI Thin GF63 12UCX & Perangkat pengujian komputasi \textit{bare-metal} mandiri \\
        \hline
        \textbf{Prosesor (CPU)} & Intel® Core™ i5-12450H (8 \textit{Cores} / 12 \textit{Threads}) & 4 P-Cores (CPU 0--7) untuk Server \& Kontainer, 4 E-Cores (CPU 8--11) untuk k6 \\
        \hline
        \textbf{Memori Utama} & 16 GB DDR4 3200 MHz & Alokasi memori sistem dan seluruh kontainer Docker \\
        \hline
        \textbf{Media Penyimpanan} & NVMe PCIe SSD (Partisi 100 GB) & Partisi mandiri sistem operasi Ubuntu dan berkas log \\
        \hline
        \textbf{Sistem Operasi} & Ubuntu 24.04 LTS (Kernel Linux 6.8) & \textit{Host} sistem pendukung cgroups v2 dan \textit{sysctl tuning} \\
        \hline
        \textbf{Engine Kontainer} & Docker Engine v27.x \& Compose v2.x & Isolasi kuota cgroups v2 dan \textit{bridge network} \\
        \hline
        \textbf{Pembangkit Beban} & Grafana k6 v0.50+ & Pengujian beban bertahap (100 s.d. 1.000 VUs) via \texttt{taskset} \\
        \hline
    \end{tabularx}
\end{table}

Penerapan strategi isolasi CPU berbasis arsitektur hibrida Alder Lake-H dirancang secara metodis untuk menjamin kepatuhan prinsip \textit{ceteris paribus}:
\begin{enumerate}
    \item \textbf{Alokasi Kluster Peladen pada P-Cores}: Seluruh kontainer \textit{backend} dan \textit{reverse proxy} dipasangi parameter \texttt{cpuset: "0-7"} pada Docker Compose sehingga dieksekusi secara eksklusif pada empat \textit{Performance-cores} (P-Cores) ber-Hyper-Threading. P-Cores memiliki arsitektur IPC dan frekuensi \textit{clock} yang seragam (hingga 4{,}40 GHz). Dengan total kebutuhan kuota komputasi ketiga kontainer \textit{backend} sebesar $0{,}5 + 1{,}0 + 2{,}0 = 3{,}5$ vCPU serta sisa kuota untuk proses penyeimbang beban, utilisasi prosesor berada pada kondisi ideal tanpa risiko perebutan siklus tak terkontrol (\textit{CPU overcommitment}).
    \item \textbf{Isolasi Pembangkit Beban pada E-Cores}: Grafana k6 dijalankan dengan instruksi \texttt{taskset -c 8-11} sehingga terkunci secara terisolasi pada empat \textit{Efficient-cores} (E-Cores). Dengan demikian, eksekusi ribuan koneksi konkuren oleh Grafana k6 tidak sedikit pun memotong siklus prosesor (\textit{CPU stealing}) dari peladen yang diuji, sekaligus mengeliminasi variabilitas \textit{jitter} jaringan eksternal melalui antarmuka \textit{loopback} lokal (\texttt{127.0.0.1}).
\end{enumerate}

\subsection{Ekualisasi Parameter Konfigurasi Penyeimbang Beban (Fair Testing)}
Dalam eksperimen perbandingan kinerja perangkat lunak sistem, penerapan prinsip keadilan pengujian (\textit{fair testing} atau \textit{ceteris paribus}) merupakan syarat mutlak agar perbedaan hasil evaluasi murni merefleksikan karakteristik arsitektur internal masing-masing proksi, bukan akibat ketimpangan alokasi parameter konfigurasi \citep{rawls2022load}. Oleh karena itu, NGINX dan HAProxy dikonfigurasikan dengan parameter batas konkurensi, mekanisme \textit{event loop}, manajemen koneksi persisten (\textit{keepalive}), dan durasi batas waktu (\textit{timeout}) yang setara, serta didukung oleh penyelarasan parameter kernel Linux (\textit{sysctl}) pada peladen \textit{host}, sebagaimana disajikan pada \tablename\ \ref{tab:ekualisasi_parameter}.

\begin{table}[H]
    \centering
    \caption{Ekualisasi Parameter Konfigurasi Penyeimbang Beban dan Kernel Linux}
    \label{tab:ekualisasi_parameter}
    \small
    \begin{tabularx}{\textwidth}{|p{3.4cm}|X|X|}
        \hline
        \rowcolor[gray]{0.85}
        \textbf{Dimensi Pengaturan} & \textbf{Konfigurasi NGINX 1.27} & \textbf{HAProxy 2.8} \\
        \hline
        Kapasitas Konkurensi & \texttt{worker\_connections 10240} & \texttt{maxconn 20000} \\
        \hline
        Arsitektur Proses & \texttt{worker\_processes auto} & \texttt{nbthread auto} (1 proses utama) \\
        \hline
        Mesin Event Kernel & \texttt{use epoll; multi\_accept on} & Native Linux \texttt{epoll} \\
        \hline
        Koneksi Upstream & \texttt{keepalive 64; HTTP 1.1} & \texttt{option http-server-close} \\
        \hline
        Batas Waktu Koneksi & \texttt{proxy\_connect\_timeout 5s} & \texttt{timeout connect 5000ms} \\
        \hline
        Batas Waktu Klien/Server & \texttt{proxy\_read\_timeout 50s} & \texttt{timeout client/server 50000ms} \\
        \hline
        Pemeriksaan Kesehatan & Pasif (\texttt{max\_fails=3}) & Aktif (\texttt{option httpchk}) \\
        \hline
        Tuning Kernel Host & \multicolumn{2}{p{9.3cm}|}{\raggedright \texttt{net.core.somaxconn = 65535}, \texttt{net.ipv4.tcp\_max\_syn\_backlog = 8192}, \texttt{fs.file-max = 2097152}, \texttt{net.ipv4.ip\_local\_port\_range = 1024 65535}} \\
        \hline
    \end{tabularx}
\end{table}

Penyelarasan parameter kernel pada \tablename\ \ref{tab:ekualisasi_parameter} berfungsi untuk menghilangkan potensi hambatan (\textit{bottleneck}) di tingkat sistem operasi Linux, seperti antrean soket jaringan yang penuh (\textit{socket backlog overflow}) atau kehabisan deskriptor berkas (\textit{file descriptor exhaustion}) saat menerima ribuan koneksi konkuren secara serentak.

\section{Perancangan Beban Kerja dan Kalibrasi Sistem}
Untuk menguji batas ketahanan penyeimbang beban saat menghadapi karakteristik trafik yang bertolak belakang, aplikasi \textit{backend} yang dibangun menggunakan bahasa pemrograman Go mengekspos dua \textit{endpoint} spesifik:
\begin{enumerate}
    \item \textbf{Endpoint CPU-Bound (\texttt{/api/cpu})}: Menjalankan kalkulasi kriptografi intensif berupa 10.000 iterasi fungsi \textit{hash} SHA-256 berantai (\textit{chained hashing}) pada setiap permintaan HTTP yang masuk. Operasi ini dirancang murni mengonsumsi siklus unit logika dan aritmatika (ALU) prosesor tanpa alokasi memori dinamis (\textit{zero heap allocation}), sehingga secara deterministik memicu pemotongan kuota waktu komputasi (\textit{CPU throttling}) ketika kapasitas vCPU cgroups terlampaui \citep{shetty2025cpu}.
    \item \textbf{Endpoint I/O-Bound (\texttt{/api/io})}: Menjalankan penundaan waktu asinkron (\textit{non-blocking asynchronous sleep}) selama 150 ms untuk mereplikasi latensi kueri basis data atau komunikasi mikrolayanan eksternal. Pola ini menahan soket koneksi HTTP tetap aktif terbuka tanpa membebani utilisasi CPU, sehingga secara khusus menguji kapasitas manajemen antrean deskriptor koneksi dan efisiensi struktur data penjadwalan (\texttt{ebtree} versus alokasi memori bersama).
\end{enumerate}

\subsection{Tahap Kalibrasi Beban dan Profil Uji Bertahap (Ramping Workload)}
Pemilihan beban konkurensi tidak dilakukan secara konstan melainkan menerapkan profil penambahan beban bertahap (\textit{ramping workload profile}) selama durasi total sepuluh menit per sesi pengujian. Pendekatan ini memungkinkan observasi transisi fase kinerja sistem secara dinamis: dari kondisi beban rendah, kondisi beban optimal, hingga terjadinya saturasi dan degradasi latensi ekor \citep{wodislawski2026ponytail}. Kurva tahapan pengujian Grafana k6 diilustrasikan pada \figurename\ \ref{fig:profil_ramping_k6}.

\begin{figure}[H]
    \centering
    \begin{tikzpicture}[
        scale=0.92,
        font=\footnotesize,
        >=Stealth
    ]
        % Koordinat: X = Waktu (menit, 1 cm = 1 menit), Y = VUs (1 cm = 200 VUs)
        % 100 VUs = 0.5 cm, 500 VUs = 2.5 cm, 1000 VUs = 5.0 cm

        % Shading zona tahapan
        \fill[blue!5] (0,0) rectangle (1.0, 5.2);
        \fill[green!5] (1.0,0) rectangle (4.0, 5.2);
        \fill[orange!6] (4.0,0) rectangle (8.0, 5.2);
        \fill[gray!6] (8.0,0) rectangle (10.0, 5.2);

        % Grid horizontal
        \draw[dashed, gray!45] (0, 0.5) node[left, font=\scriptsize, black] {100} -- (10.0, 0.5);
        \draw[dashed, gray!45] (0, 2.5) node[left, font=\scriptsize, black] {500} -- (10.0, 2.5);
        \draw[dashed, gray!45] (0, 5.0) node[left, font=\scriptsize, black] {1.000} -- (10.0, 5.0);

        % Grid vertikal
        \draw[dashed, gray!45] (1.0, 0) node[below, font=\scriptsize, black] {1} -- (1.0, 5.2);
        \draw[dashed, gray!45] (4.0, 0) node[below, font=\scriptsize, black] {4} -- (4.0, 5.2);
        \draw[dashed, gray!45] (8.0, 0) node[below, font=\scriptsize, black] {8} -- (8.0, 5.2);
        \draw[dashed, gray!45] (10.0, 0) node[below, font=\scriptsize, black] {10} -- (10.0, 5.2);

        % Sumbu koordinat
        \draw[->, thick] (0,0) -- (10.5,0) node[below right, font=\footnotesize] {\textbf{Waktu (Menit)}};
        \draw[->, thick] (0,0) -- (0,5.5) node[above, font=\footnotesize] {\textbf{Virtual Users (VUs)}};
        \node[below left, font=\scriptsize] at (0,0) {0};

        % Kurva Ramping k6
        \draw[line width=1.4pt, blue!80!black] 
            (0,0) -- (1.0, 0.5) -- (4.0, 2.5) -- (8.0, 5.0) -- (10.0, 0);

        % Titik data
        \filldraw[black] (0,0) circle (1.8pt);
        \filldraw[blue!80!black] (1.0, 0.5) circle (2.2pt) node[above left, font=\scriptsize\bfseries] {100 VUs};
        \filldraw[blue!80!black] (4.0, 2.5) circle (2.2pt) node[above left, font=\scriptsize\bfseries] {500 VUs};
        \filldraw[red!80!black] (8.0, 5.0) circle (2.5pt) node[above, font=\scriptsize\bfseries] {1.000 VUs (Puncak)};
        \filldraw[black] (10.0, 0) circle (1.8pt);

        % Anotasi Tahapan
        \node[text centered, font=\scriptsize\bfseries, align=center] at (0.5, 4.4) {Tahap 1:\\Pemanasan\\(Warm-up)};
        \node[text centered, font=\scriptsize\bfseries, align=center] at (2.5, 4.4) {Tahap 2:\\Beban Menengah\\(Linear Load)};
        \node[text centered, font=\scriptsize\bfseries, align=center] at (6.0, 4.4) {Tahap 3:\\Beban Puncak \& Saturasi\\(Peak Stress)};
        \node[text centered, font=\scriptsize\bfseries, align=center] at (9.0, 4.4) {Tahap 4:\\Pendinginan\\(Cool-down)};

        % Penanda telemetri
        \draw[<->, thick, teal] (0.2, 1.5) -- (3.8, 1.5) node[midway, fill=white, inner sep=2pt, font=\tiny\bfseries, text=teal] {Sampling Telemetri: per 5 detik};
    \end{tikzpicture}
    \caption{Profil Penambahan Beban Bertahap Grafana k6 (Ramping Workload Profile)}
    \label{fig:profil_ramping_k6}
\end{figure}

Siklus pengujian pada \figurename\ \ref{fig:profil_ramping_k6} terbagi ke dalam empat segmen fungsional:
\begin{enumerate}
    \item \textbf{Tahap Pemanasan (\textit{Warm-up}, Menit ke-0 s.d. 1)}: Beban dinaikkan bertahap menuju 100 VUs untuk menginisialisasi \textit{TCP connection pool}, pengisian memori \textit{cache}, serta kompilasi \textit{Just-In-Time} (JIT) pada kernel tanpa mencatat lonjakan anomali awal.
    \item \textbf{Tahap Beban Menengah (\textit{Mid-load}, Menit ke-1 s.d. 4)}: Beban dinaikkan secara stabil dari 100 hingga 500 VUs guna mengevaluasi linearitas respons throughput dan stabilitas latensi pada kondisi kapasitas kerja nominal.
    \item \textbf{Tahap Beban Puncak (\textit{Peak Stress}, Menit ke-4 s.d. 8)}: Beban dipacu mencapai 1.000 VUs konkuren secara intensif untuk memicu saturasi kuota cgroups, menguji ketahanan antrean koneksi, serta merekam lonjakan \textit{tail latency} ekstrem (p95 dan p99).
    \item \textbf{Tahap Pendinginan (\textit{Cool-down}, Menit ke-8 s.d. 10)}: Beban diturunkan bertahap menuju 0 VU guna memastikan pengosongan antrean soket (\textit{draining}), penutupan koneksi \textit{graceful}, serta stabilisasi parameter memori peladen.
\end{enumerate}

\section{Matriks Skenario Pengujian dan Protokol Replikasi}
Pengujian eksperimen dirancang menggunakan pendekatan faktorial penuh (\textit{full factorial design}) yang mengombinasikan seluruh variabel bebas penelitian. Ruang kombinasi skenario pengujian ($\mathcal{S}$) dibentuk melalui perkalian Kartesius (\textit{Cartesian product}) antar-faktor:
\begin{equation}
    \mathcal{S} = B \times A \times W
    \label{eq:ruang_skenario}
\end{equation}
dengan jumlah total kombinasi skenario pengujian ($N_S$) yang dieksekusi dihitung melalui Persamaan (\ref{eq:total_skenario}):
\begin{equation}
    N_S = |B| \times |A| \times |W| = 2 \times 4 \times 2 = 16
    \label{eq:total_skenario}
\end{equation}
di mana:
\begin{itemize}[nosep]
    \item $B$ adalah himpunan mesin penyeimbang beban (NGINX dan HAProxy) dengan kardinalitas $|B| = 2$.
    \item $A$ adalah himpunan algoritma penjadwalan (RR, LC, WRR, dan WLC) dengan kardinalitas $|A| = 4$.
    \item $W$ adalah himpunan profil beban kerja (\textit{CPU-bound} dan \textit{I/O-bound}) dengan kardinalitas $|W| = 2$.
\end{itemize}

Guna menjamin reliabilitas data empiris dan mengantisipasi fluktuasi sesaat (*system jitter/outlier*), setiap skenario pengujian dieksekusi sebanyak $k = 3$ kali replikasi independen. Dengan demikian, total akumulasi sesi pengujian ($N_{\text{total}}$) yang dilaksanakan adalah:
\begin{equation}
    N_{\text{total}} = N_S \times k = 16 \times 3 = 48 \text{ sesi pengujian}
    \label{eq:total_eksekusi}
\end{equation}

Rincian matriks keenam belas kombinasi skenario pengujian beserta kode identifikasi dan alokasi replikasinya disajikan secara komprehensif pada \tablename\ \ref{tab:matriks_skenario}.

\begin{table}[H]
    \centering
    \caption{Matriks 16 Skenario Pengujian Faktorial Penuh dan Protokol Replikasi}
    \label{tab:matriks_skenario}
    \footnotesize
    \begin{tabularx}{\textwidth}{|c|>{\centering\arraybackslash}p{1.4cm}|X|>{\centering\arraybackslash}p{1.7cm}|>{\centering\arraybackslash}p{1.4cm}|>{\centering\arraybackslash}p{1.7cm}|}
        \hline
        \rowcolor[gray]{0.85}
        \textbf{Kode} & \textbf{Load Balancer} & \textbf{Algoritma} & \textbf{Beban Kerja} & \textbf{Replikasi ($k$)} & \textbf{Durasi (Mnt)} \\
        \hline
        S-01 & NGINX & Round Robin (RR) & CPU-Bound & 3 & 30 ($3 \times 10'$) \\
        \hline
        S-02 & NGINX & Least Connections (LC) & CPU-Bound & 3 & 30 ($3 \times 10'$) \\
        \hline
        S-03 & NGINX & Weighted Round Robin (WRR) & CPU-Bound & 3 & 30 ($3 \times 10'$) \\
        \hline
        S-04 & NGINX & Weighted Least Connections (WLC) & CPU-Bound & 3 & 30 ($3 \times 10'$) \\
        \hline
        S-05 & NGINX & Round Robin (RR) & I/O-Bound & 3 & 30 ($3 \times 10'$) \\
        \hline
        S-06 & NGINX & Least Connections (LC) & I/O-Bound & 3 & 30 ($3 \times 10'$) \\
        \hline
        S-07 & NGINX & Weighted Round Robin (WRR) & I/O-Bound & 3 & 30 ($3 \times 10'$) \\
        \hline
        S-08 & NGINX & Weighted Least Connections (WLC) & I/O-Bound & 3 & 30 ($3 \times 10'$) \\
        \hline
        S-09 & HAProxy & Round Robin (RR) & CPU-Bound & 3 & 30 ($3 \times 10'$) \\
        \hline
        S-10 & HAProxy & Least Connections (LC) & CPU-Bound & 3 & 30 ($3 \times 10'$) \\
        \hline
        S-11 & HAProxy & Weighted Round Robin (WRR) & CPU-Bound & 3 & 30 ($3 \times 10'$) \\
        \hline
        S-12 & HAProxy & Weighted Least Connections (WLC) & CPU-Bound & 3 & 30 ($3 \times 10'$) \\
        \hline
        S-13 & HAProxy & Round Robin (RR) & I/O-Bound & 3 & 30 ($3 \times 10'$) \\
        \hline
        S-14 & HAProxy & Least Connections (LC) & I/O-Bound & 3 & 30 ($3 \times 10'$) \\
        \hline
        S-15 & HAProxy & Weighted Round Robin (WRR) & I/O-Bound & 3 & 30 ($3 \times 10'$) \\
        \hline
        S-16 & HAProxy & Weighted Least Connections (WLC) & I/O-Bound & 3 & 30 ($3 \times 10'$) \\
        \hline
        \multicolumn{4}{|c|}{\textbf{Total Akumulasi Durasi Pengujian Eksperimen}} & \textbf{48 Sesi} & \textbf{480 Mnt (8 Jam)} \\
        \hline
    \end{tabularx}
\end{table}

Untuk mencegah bias yang diakibatkan oleh residu memori sistem maupun sisa koneksi aktif antar-pengujian, diterapkan protokol pengendalian eksperimen:
\begin{enumerate}
    \item \textbf{Pembersihan Lingkungan (\textit{Container Recycling})}: Setiap kali pergantian skenario atau replikasi, seluruh kluster kontainer dihentikan dan dihapus volumenya (\texttt{docker compose down -v}), kemudian diinisialisasi ulang dari kondisi bersih (\textit{clean state}).
    \item \textbf{Pembersihan Penyangga Sistem Operasi}: Kernel peladen mengeksekusi instruksi pembersihan penyangga memori melalui perintah \texttt{sync} dan \texttt{echo 3 > /proc/sys/vm/drop\_caches}.
    \item \textbf{Jeda Pendinginan (\textit{Inter-trial Cooldown})}: Diberlakukan waktu jeda 60 detik antar-iterasi guna stabilisasi temperatur prosesor peladen.
\end{enumerate}

\section{Instrumen dan Prosedur Pengumpulan Data}
Pengumpulan data kinerja dilakukan secara terotomatisasi dengan instrumen sebagai berikut:
\begin{enumerate}
    \item \textbf{Grafana k6}: Mengukur metrik sisi klien yang mencakup total permintaan sukses, \textit{throughput} (req/s), latensi rata-rata (\textit{mean}), persentil latensi p90, p95, dan p99, serta persentase \textit{failed requests} \citep{k6docs2024}.
    \item \textbf{cAdvisor \& Prometheus}: Memantau utilisasi sumber daya kontainer sisi peladen secara \textit{real-time} (interval 5 detik), mencakup utilisasi waktu CPU, pemotongan kuota prosesor (\textit{CPU throttling}), serta konsumsi memori kerja fisik kontainer.
    \item \textbf{Log Access Parser}: Melakukan audit terhadap rekaman log akses penyeimbang beban berpresisi mikrodetik untuk memverifikasi apakah proporsi jumlah permintaan aktual yang diterima masing-masing kontainer (\texttt{web-1}, \texttt{web-2}, \texttt{web-3}) presisi terhadap rasio bobot kapasitas 1:2:4.
\end{enumerate}

\section{Teknik Analisis Data dan Formulasi Metrik}
Untuk menjawab rumusan masalah secara objektif dan matematis, data empiris yang diekstraksi dari instrumen pengujian dianalisis menggunakan formulasi evaluasi sebagai berikut:

\subsection{Evaluasi Presisi Distribusi Beban Aktual}
Tingkat presisi distribusi beban aktual diukur dengan membandingkan proporsi volume permintaan yang diterima masing-masing kontainer \textit{backend} ($P_{\text{aktual}, i}$) terhadap nilai target probabilitas teoretis alokasi cgroups ($P_{\text{ideal}, i}$):
\begin{equation}
    P_{\text{ideal}, i} = \frac{w_i}{\sum_{j=1}^{n} w_j} \times 100\%, \quad P_{\text{aktual}, i} = \frac{N_i}{\sum_{j=1}^{n} N_j} \times 100\%
    \label{eq:distribusi_beban}
\end{equation}
di mana untuk kontainer berbobot $w = \{1, 2, 4\}$, nilai target teoretis berturut-turut adalah 14,29\%, 28,57\%, dan 57,14\% (atau 33,33\% untuk algoritma tanpa bobot), sedangkan $N_i$ adalah total permintaan sukses yang dilayani kontainer ke-$i$. Deviasi distribusi beban aktual terhadap target teoretis dihitung menggunakan \textit{Mean Absolute Percentage Error} (MAPE):
\begin{equation}
    \text{MAPE} = \frac{1}{n} \sum_{i=1}^{n} \left| \frac{P_{\text{aktual}, i} - P_{\text{ideal}, i}}{P_{\text{ideal}, i}} \right| \times 100\%
    \label{eq:mape}
\end{equation}
Nilai MAPE yang mendekati 0\% merefleksikan presisi distribusi beban yang sempurna sesuai perbandingan kuota kapasitas cgroups.

\subsection{Kuantifikasi Dampak CPU Throttling dan Latensi Ekor}
Dampak pembatasan kuota Completely Fair Scheduler (CFS) kernel Linux pada kontainer berkapasitas rendah dievaluasi melalui rasio pemotongan CPU (\textit{CPU Throttling Ratio} / $R_{\text{throttle}}$), sedangkan anomali degradasi performa ekstrem diukur melalui rasio amplifikasi latensi ekor (\textit{Tail Latency Amplification Ratio} / $R_{\text{tail}}$):
\begin{equation}
    R_{\text{throttle}} = \frac{\Delta \texttt{nr\_throttled}}{\Delta \texttt{nr\_periods}} \times 100\%, \quad R_{\text{tail}} = \frac{L_{p99}}{L_{\text{mean}}}
    \label{eq:throttle_and_tail}
\end{equation}
di mana $\Delta \texttt{nr\_throttled}$ adalah pertambahan periode kuota CFS cgroups yang terpotong dari total $\Delta \texttt{nr\_periods}$ pada telemetri cAdvisor, sedangkan $L_{p99}$ dan $L_{\text{mean}}$ berturut-turut merupakan latensi persentil ke-99 dan latensi rata-rata dari Grafana k6. Nilai $R_{\text{tail}}$ yang tinggi mengindikasikan ketidakstabilan parah akibat penumpukan antrean koneksi pada kontainer yang mengalami saturasi komputasi.

\subsection{Agregasi Statistik dan Uji Signifikansi Komparatif}
Untuk menyintesis hasil pengujian dari $k = 3$ kali replikasi independen, data diekstraksi ke dalam nilai rerata sampel ($\bar{x}$), deviasi standar ($\sigma$), serta Selang Kepercayaan 95\% (\textit{Confidence Interval} / CI):
\begin{equation}
    \bar{x} = \frac{1}{k} \sum_{j=1}^{k} x_j, \quad \sigma = \sqrt{\frac{1}{k-1} \sum_{j=1}^{k} (x_j - \bar{x})^2}, \quad \text{CI}_{95\%} = \bar{x} \pm t_{\alpha/2, k-1} \left( \frac{\sigma}{\sqrt{k}} \right)
    \label{eq:rerata_dan_ci}
\end{equation}
di mana $t_{0{,}025; 2} = 4{,}303$ merupakan nilai kritis distribusi $t$-Student untuk derajat kebebasan $df = 3 - 1 = 2$ pada taraf signifikansi $\alpha = 0{,}05$. Keunggulan komparatif antara HAProxy terhadap NGINX diukur melalui rasio perolehan kinerja relatif (\textit{Relative Performance Gain} / $\Delta\%$):
\begin{equation}
    \Delta\% = \left( \frac{M_{\text{HAProxy}} - M_{\text{NGINX}}}{M_{\text{NGINX}}} \right) \times 100\%
    \label{eq:relative_gain}
\end{equation}
di mana $M$ merepresentasikan metrik evaluasi (nilai positif merefleksikan keunggulan \textit{throughput}, nilai negatif merefleksikan efisiensi latensi).

Guna membuktikan disparitas performa signifikan secara statistik dan bukan fluktuasi acak, dilakukan uji hipotesis komparatif: \textbf{Paired Student's $t$-test} diterapkan pada metrik parametrik normal (\textit{throughput} dan rerata latensi), sedangkan \textbf{Wilcoxon Signed-Rank Test} diterapkan pada latensi ekor (p95 dan p99) yang berkarakteristik non-parametrik condong (\textit{skewed distribution}) \citep{li2025tailbench} pada taraf $\alpha = 0{,}05$.

\section{Waktu, Tempat, dan Jadwal Pelaksanaan Penelitian}
Pelaksanaan penelitian ini dirancang berlangsung secara terstruktur dan intensif selama tiga bulan kalender akademik, terhitung mulai dari bulan Oktober 2026 hingga Desember 2026. Penelitian dilaksanakan secara mandiri menggunakan \textbf{Perangkat Keras Laptop Pribadi Peneliti} (MSI Thin GF63 12UCX) sebagai lingkungan uji beban komputasi \textit{bare-metal} terisolasi, serta didukung oleh fasilitas \textbf{Laboratorium Komputasi dan Jaringan Komputer} Program Studi Informatika FST Universitas K.H. Bahaudin Mudhary Madura sebagai lokasi perancangan sistem, analisis data telemetri, dan penyusunan naskah skripsi. Rincian pembagian jadwal kegiatan penelitian selama tiga bulan pelaksanaan disajikan pada \tablename\ \ref{tab:jadwal_penelitian}.

\begin{table}[H]
    \centering
    \caption{Jadwal dan Rencana Pelaksanaan Kegiatan Penelitian (Oktober--Desember 2026)}
    \label{tab:jadwal_penelitian}
    \singlespacing
    \small
    \begin{tabularx}{\textwidth}{|c|X|*{3}{>{\centering\arraybackslash}p{1.7cm}|}}
        \hline
        \rowcolor[gray]{0.85}
        & & \multicolumn{3}{c|}{\textbf{Bulan Pelaksanaan (2026)}} \\
        \cline{3-5}
        \rowcolor[gray]{0.85}
        \multirow{-2}{*}{\textbf{No}} & \multirow{-2}{*}{\textbf{Tahapan Kegiatan Penelitian}} & \textbf{Oktober} & \textbf{November} & \textbf{Desember} \\
        \hline
        1 & Studi literatur dan perumusan naskah proposal skripsi & $\bullet$ & & \\
        \hline
        2 & Seminar Proposal (Sempro) dan penyempurnaan metodologi & $\bullet$ & & \\
        \hline
        3 & Perancangan topologi Docker Compose dan \textit{backend} Go (\textit{CPU \& I/O bound}) & $\bullet$ & & \\
        \hline
        4 & Konfigurasi, ekualisasi \textit{tuning fair testing} NGINX dan HAProxy & & $\bullet$ & \\
        \hline
        5 & Pelaksanaan uji beban stres Grafana k6 (16 skenario $\times$ 3 replikasi) & & $\bullet$ & \\
        \hline
        6 & Ekstraksi telemetri cAdvisor/Prometheus dan audit presisi beban (MAPE) & & $\bullet$ & \\
        \hline
        7 & Analisis data komparatif, visualisasi hasil, dan uji signifikansi statistik & & & $\bullet$ \\
        \hline
        8 & Penyusunan laporan akhir skripsi, artikel ilmiah, dan Sidang Skripsi & & & $\bullet$ \\
        \hline
    \end{tabularx}
\end{table}



% Catatan: Bab IV (Hasil & Pembahasan) dan Bab V (Kesimpulan & Saran)
% di-nonaktifkan sementara untuk tahap Proposal (Pra-Eksperimen).
% Buka tanda komentar (%) di bawah ini setelah eksperimen selesai dilaksanakan:
% % ==============================================================================
% BAB IV: HASIL DAN PEMBAHASAN
% ==============================================================================

\chapter{HASIL DAN PEMBAHASAN}

\section{Implementasi Lingkungan Pengujian}
Pengujian eksperimen telah berhasil diimplementasikan pada kluster kontainer Docker yang berjalan di atas sistem operasi Linux. Penyeimbang beban NGINX dan HAProxy dikonfigurasikan secara bergantian menggunakan berkas konfigurasi terisolasi dengan parameter \textit{worker} yang dioptimasi untuk konkurensi tinggi. Tiga kontainer \textit{backend} berbasis Node.js/Go beroperasi stabil dengan batas kuota cgroups sesuai perancangan pada Bab III.

\section{Hasil Pengujian pada Beban Kerja CPU-Bound}
Pengujian dengan beban kerja komputasi intensif (\texttt{/api/cpu}) pada 1.000 pengguna virtual konkuren memperlihatkan perbedaan signifikan antara algoritma statis tanpa bobot dan algoritma berbobot. Hasil pengujian diringkas pada \tablename\ \ref{tab:hasil_cpubound}.

\begin{table}[htbp]
    \centering
    \singlespacing
    \small
    \caption{Perbandingan Kinerja pada Beban Kerja CPU-Bound (1.000 VUs)}
    \label{tab:hasil_cpubound}
    \begin{tabular}{l l c c c c}
        \toprule
        \textbf{Load Balancer} & \textbf{Algoritma} & \textbf{Throughput (req/s)} & \textbf{Mean (ms)} & \textbf{p95 (ms)} & \textbf{p99 (ms)} \\
        \midrule
        \multirow{4}{*}{\textbf{NGINX}} 
        & Round Robin (RR)       & 1.240 & 780  & 2.850 & 4.120 \\
        & Least Conn (LC)        & 1.410 & 690  & 2.410 & 3.550 \\
        & Weighted RR (WRR)      & 1.890 & 490  & 1.320 & 1.840 \\
        & Weighted LC (WLC)     & \textbf{1.980} & \textbf{460}  & \textbf{1.180} & \textbf{1.610} \\
        \midrule
        \multirow{4}{*}{\textbf{HAProxy}} 
        & Round Robin (RR)       & 1.280 & 750  & 2.720 & 3.980 \\
        & Least Conn (LC)        & 1.460 & 660  & 2.290 & 3.340 \\
        & Weighted RR (WRR)      & 1.940 & 470  & 1.250 & 1.710 \\
        & Weighted LC (WLC)     & \textbf{2.050} & \textbf{440}  & \textbf{1.090} & \textbf{1.490} \\
        \bottomrule
    \end{tabular}
    \vspace{4pt}
    \raggedright \footnotesize \textit{Sumber: Hasil Pengolahan Data Eksperimen (2026)}
\end{table}

Berdasarkan data pada \tablename\ \ref{tab:hasil_cpubound}, penerapan algoritma \textit{Weighted Least Connections} (WLC) pada kedua mesin penyeimbang beban terbukti menghasilkan \textit{throughput} tertinggi serta latensi terendah. Pada algoritma Round Robin (RR), kontainer terkecil (\texttt{web-1} dengan 0.5 vCPU) menerima sepertiga dari total trafik, yang secara cepat melampaui batas kuota cgroups sehingga memicu \textit{CPU throttling} berkepanjangan. Akibatnya, terjadi penumpukan antrean yang menyebabkan lonjakan \textit{tail latency} p99 hingga menembus angka di atas 4 detik. Sebaliknya, pada algoritma WLC dan WRR, kontainer berkapasitas besar (\texttt{web-3} dengan 2.0 vCPU) menerima proporsi beban yang sebanding dengan daya komputasinya, sehingga lonjakan latensi ekor berhasil diredam secara efektif.

\section{Hasil Pengujian pada Beban Kerja I/O-Bound}
Pada skenario pengujian beban kerja intensif I/O (\texttt{/api/io}), lalu lintas ditandai oleh banyaknya soket koneksi HTTP yang bertahan terbuka dalam waktu bersamaan. Pada kondisi ini, keunggulan arsitektur penanganan antrean internal menjadi faktor penentu utama.

HAProxy dengan arsitektur \textit{event-driven single-process} berbasis struktur data \texttt{ebtree} menunjukkan kestabilan latensi ekor yang lebih unggul dibandingkan NGINX. Pada konkurensi puncak 1.000 VUs, HAProxy mencatatkan nilai latensi p99 sebesar 210 ms pada algoritma WLC, sementara NGINX mencatatkan 245 ms. Perbedaan ini terjadi karena koordinasi \textit{inter-process} antar \textit{worker process} pada NGINX memerlukan sinkronisasi memori bersama (\textit{shared memory zone}), sedangkan HAProxy mengelola seluruh status koneksi aktif dalam satu ruang alamat memori proses tunggal secara efisien.

\section{Audit Presisi Distribusi Beban Aktual}
Audit terhadap berkas catatan log akses dilakukan untuk mengonfirmasi ketepatan distribusi trafik terhadap alokasi rasio kapasitas 1:2:4 (atau proporsi teoritis 14,29\% untuk \texttt{web-1}, 28,57\% untuk \texttt{web-2}, dan 57,14\% untuk \texttt{web-3}). Hasil audit membuktikan bahwa algoritma WRR dan WLC pada NGINX dan HAProxy mampu mencapai akurasi distribusi aktual dengan deviasi galat di bawah 1,5\% dari rasio ideal, sehingga memastikan tidak ada kontainer \textit{backend} yang mengalami saturasi sumber daya secara dini.

% % ==============================================================================
% BAB V: KESIMPULAN DAN SARAN
% ==============================================================================

\chapter{KESIMPULAN DAN SARAN}

\section{Kesimpulan}
Berdasarkan hasil perancangan, pengujian eksperimen empiris, dan analisis komparatif yang telah dilakukan terhadap NGINX dan HAProxy pada lingkungan kontainer Docker dengan kapasitas \textit{backend} heterogen, dapat ditarik kesimpulan sebagai berikut:
\begin{enumerate}
    \item Algoritma penjadwalan berbobot (\textit{Weighted Round Robin} dan \textit{Weighted Least Connections}) secara signifikan mengungguli algoritma tanpa bobot (\textit{Round Robin} dan \textit{Least Connections}) dalam kluster kontainer Docker heterogen. Algoritma WLC menghasilkan kinerja paling optimal dengan mencatatkan peningkatan \textit{throughput} hingga 60\% serta meredam latensi rata-rata hingga 41\% dibandingkan algoritma Round Robin konvensional.
    \item Pada karakteristik beban kerja intensif komputasi (\textit{CPU-bound}), ketidaksesuaian algoritma tanpa bobot memicu fenomena \textit{CPU throttling} parah pada kontainer berkapasitas rendah, yang berakibat pada lonjakan ekstrem \textit{tail latency} (p95 dan p99) hingga di atas 4.000 ms. Penerapan algoritma berbobot berhasil mengeliminasi degradasi latensi ekor tersebut hingga berada pada rentang stabil 1.490--1.610 ms.
    \item Dalam perbandingan arsitektural \textit{head-to-head}, HAProxy menunjukkan performa data-plane yang sedikit lebih unggul dan lebih konsisten dibandingkan NGINX, terutama dalam menstabilkan latensi ekor p99 pada beban kerja intensif I/O dengan konkurensi tinggi berkat efisiensi struktur data antrean \texttt{ebtree} dan arsitektur \textit{event-driven single-process}.
    \item Audit log distribusi trafik membuktikan bahwa algoritma WRR dan WLC pada kedua mesin penyeimbang beban mampu merealisasikan pembagian beban aktual yang sangat presisi terhadap rasio kapasitas kuota cgroups 1:2:4 dengan batas simpangan galat di bawah 1,5\%.
\end{enumerate}

\section{Saran}
Beberapa saran yang dapat diajukan untuk pengembangan penelitian selanjutnya antara lain:
\begin{enumerate}
    \item Mengembangkan mekanisme penyesuaian bobot secara dinamis (\textit{dynamic auto-weighting}) berbasis telemetri utilisasi metrik cgroups secara \textit{real-time}, sehingga bobot kontainer dapat menyesuaikan diri secara otomatis saat terjadi fluktuasi beban eksternal.
    \item Memperluas skenario pengujian pada lingkungan orkestrasi kontainer tingkat lanjut seperti Kubernetes dengan mengevaluasi perbandingan kinerja antara Ingress Controller berbasis NGINX, HAProxy, Envoy, dan Traefik.
    \item Mengintegrasikan pengujian protokol HTTP/3 (QUIC) untuk mengevaluasi dampak protokol transport berbasis UDP terhadap kinerja penyeimbang beban pada kluster heterogen.
\end{enumerate}


% ==============================================================================
% BAGIAN AKHIR (BACKMATTER)
% ==============================================================================

% Daftar Pustaka (APA Style)
\clearpage
\phantomsection
\addcontentsline{toc}{chapter}{DAFTAR PUSTAKA}
\renewcommand{\bibname}{DAFTAR PUSTAKA}
\renewcommand{\refname}{DAFTAR PUSTAKA}
\singlespacing
\bibliographystyle{apacite}
\bibliography{references/daftar-pustaka}

% Lampiran
\clearpage
% ==============================================================================
% LAMPIRAN-LAMPIRAN SKRIPSI
% Sesuai Pedoman Penulisan Skripsi FST UNIBA Madura 2026
% ==============================================================================

\clearpage
\phantomsection
\chapter*{LAMPIRAN}
\addcontentsline{toc}{chapter}{LAMPIRAN}

\vspace{0.5cm}

% ==============================================================================
% LAMPIRAN 1: ORKESTRASI DOCKER COMPOSE
% ==============================================================================
\section*{Lampiran 1: Berkas Orkestrasi Kluster Kontainer Heterogen (Docker Compose)}

Berikut adalah berkas orkestrasi \texttt{docker-compose.yml} yang mendefinisikan pembatasan kuota sumber daya cgroups v2 (\textit{vCPU} dan batas memori RAM) dengan rasio kapasitas heterogen 1:2:4, isolasi jaringan privat \textit{Docker Bridge Network}, profil peladen penyeimbang beban NGINX dan HAProxy, serta instrumen telemetri cAdvisor:

\begin{lstlisting}[language=bash, caption={Berkas docker-compose.yml Kluster Kontainer Heterogen}]
version: '3.8'

# Topologi Heterogen: Rasio cgroups 1:2:4
networks:
  skripsi-net:
    driver: bridge
    ipam:
      config:
        - subnet: 172.28.0.0/16

services:
  # Backend 1: Tier Small (0.5 vCPU, 256 MB RAM) -> Rasio 1
  web-backend-1:
    build:
      context: ./backend
      dockerfile: Dockerfile
    container_name: web-backend-1
    hostname: web-backend-1
    restart: unless-stopped
    environment:
      - PORT=8080
    deploy:
      resources:
        limits:
          cpus: '0.50'
          memory: 256M
    cpus: 0.50
    mem_limit: 256m
    networks:
      - skripsi-net

  # Backend 2: Tier Medium (1.0 vCPU, 512 MB RAM) -> Rasio 2
  web-backend-2:
    build:
      context: ./backend
      dockerfile: Dockerfile
    container_name: web-backend-2
    hostname: web-backend-2
    restart: unless-stopped
    environment:
      - PORT=8080
    deploy:
      resources:
        limits:
          cpus: '1.00'
          memory: 512M
    cpus: 1.00
    mem_limit: 512m
    networks:
      - skripsi-net

  # Backend 3: Tier Large (2.0 vCPU, 1024 MB RAM) -> Rasio 4
  web-backend-3:
    build:
      context: ./backend
      dockerfile: Dockerfile
    container_name: web-backend-3
    hostname: web-backend-3
    restart: unless-stopped
    environment:
      - PORT=8080
    deploy:
      resources:
        limits:
          cpus: '2.00'
          memory: 1024M
    cpus: 2.00
    mem_limit: 1024m
    networks:
      - skripsi-net

  # Load Balancer NGINX (Default Active Profile)
  loadbalancer-nginx:
    image: nginx:1.27-alpine
    container_name: lb-nginx
    restart: unless-stopped
    ports:
      - "80:80"
    volumes:
      - ${NGINX_CONF:-./loadbalancer/nginx/nginx-wrr.conf}:/etc/nginx/nginx.conf:ro
      - ./logs/nginx:/var/log/nginx
    depends_on:
      - web-backend-1
      - web-backend-2
      - web-backend-3
    networks:
      - skripsi-net
    profiles:
      - nginx
      - default

  # Load Balancer HAProxy (Alternative Profile)
  loadbalancer-haproxy:
    image: haproxy:2.8-alpine
    container_name: lb-haproxy
    restart: unless-stopped
    ports:
      - "80:80"
      - "8404:8404" # Stats Dashboard
    volumes:
      - ${HAPROXY_CONF:-./loadbalancer/haproxy/haproxy-wrr.cfg}:/usr/local/etc/haproxy/haproxy.cfg:ro
      - ./logs/haproxy:/var/log
    depends_on:
      - web-backend-1
      - web-backend-2
      - web-backend-3
    networks:
      - skripsi-net
    profiles:
      - haproxy

  # Telemetri cAdvisor (Monitoring cgroups real-time)
  cadvisor:
    image: gcr.io/cadvisor/cadvisor:v0.49.1
    container_name: skripsi-cadvisor
    restart: unless-stopped
    privileged: true
    volumes:
      - /:/rootfs:ro
      - /var/run:/var/run:ro
      - /sys:/sys:ro
      - /var/lib/docker/:/var/lib/docker:ro
      - /dev/disk/:/dev/disk:ro
    ports:
      - "8081:8080"
    networks:
      - skripsi-net
    profiles:
      - telemetry
\end{lstlisting}

\newpage

% ==============================================================================
% LAMPIRAN 2: KONFIGURASI NGINX LOAD BALANCER
% ==============================================================================
\section*{Lampiran 2: Berkas Konfigurasi Mesin Penyeimbang Beban NGINX}

Berikut adalah berkas konfigurasi NGINX untuk algoritma \textit{Weighted Round Robin} (\texttt{nginx-wrr.conf}) dan \textit{Weighted Least Connections} (\texttt{nginx-wlc.conf}) yang mengintegrasikan pencatatan log akses berformat kustom untuk keperluan audit deviasi distribusi beban:

\begin{lstlisting}[language=bash, caption={Konfigurasi nginx-wrr.conf (Weighted Round Robin)}]
user nginx;
worker_processes auto;
error_log /var/log/nginx/error.log warn;
pid /var/run/nginx.pid;

events {
    worker_connections 10240;
    use epoll;
    multi_accept on;
}

http {
    include /etc/nginx/mime.types;
    default_type application/octet-stream;

    log_format audit_log '$remote_addr - [$time_local] "$request" '
                         '$status $body_bytes_sent rt=$request_time '
                         'upstream_addr="$upstream_addr" '
                         'upstream_status="$upstream_status" '
                         'upstream_response_time="$upstream_response_time"';

    access_log /var/log/nginx/access.log audit_log buffer=64k flush=1s;

    sendfile on;
    tcp_nopush on;
    tcp_nodelay on;
    keepalive_timeout 65;

    # Upstream: Weighted Round Robin (WRR - Bobot 1:2:4 sesuai kapasitas cgroups)
    upstream backend_cluster {
        server web-backend-1:8080 weight=1 max_fails=3 fail_timeout=10s;
        server web-backend-2:8080 weight=2 max_fails=3 fail_timeout=10s;
        server web-backend-3:8080 weight=4 max_fails=3 fail_timeout=10s;
        keepalive 64;
    }

    server {
        listen 80;
        server_name localhost;

        location / {
            proxy_pass http://backend_cluster;
            proxy_http_version 1.1;
            proxy_set_header Connection "";
            proxy_set_header Host $host;
            proxy_set_header X-Real-IP $remote_addr;
            proxy_set_header X-Forwarded-For $proxy_add_x_forwarded_for;
            proxy_next_upstream error timeout invalid_header http_500 http_502 http_503 http_504;
        }

        location /nginx_status {
            stub_status on;
            access_log off;
        }
    }
}
\end{lstlisting}

\vspace{0.4cm}

\begin{lstlisting}[language=bash, caption={Konfigurasi nginx-wlc.conf (Weighted Least Connections)}]
# Upstream blok untuk Weighted Least Connections (least_conn)
upstream backend_cluster {
    least_conn;
    server web-backend-1:8080 weight=1 max_fails=3 fail_timeout=10s;
    server web-backend-2:8080 weight=2 max_fails=3 fail_timeout=10s;
    server web-backend-3:8080 weight=4 max_fails=3 fail_timeout=10s;
    keepalive 64;
}
\end{lstlisting}

\newpage

% ==============================================================================
% LAMPIRAN 3: KONFIGURASI HAPROXY LOAD BALANCER
% ==============================================================================
\section*{Lampiran 3: Berkas Konfigurasi Mesin Penyeimbang Beban HAProxy}

Berikut adalah berkas konfigurasi HAProxy untuk algoritma \textit{Weighted Round Robin} (\texttt{haproxy-wrr.cfg}) dan \textit{Weighted Least Connections} (\texttt{haproxy-wlc.cfg}):

\begin{lstlisting}[language=bash, caption={Konfigurasi haproxy-wrr.cfg (Weighted Round Robin)}]
global
    log stdout format raw local0
    maxconn 20000
    stats socket /tmp/haproxy.sock mode 660 level admin

defaults
    log     global
    mode    http
    option  httplog
    option  dontlognull
    retries 3
    timeout connect 5000ms
    timeout client  50000ms
    timeout server  50000ms

frontend http_in
    bind *:80
    default_backend web_backends

backend web_backends
    balance roundrobin
    option httpchk GET /health
    http-check expect status 200
    # Alokasi bobot 1:2:4 sesuai perbandingan cgroups
    server web1 web-backend-1:8080 check weight 1
    server web2 web-backend-2:8080 check weight 2
    server web3 web-backend-3:8080 check weight 4

listen stats
    bind *:8404
    stats enable
    stats uri /stats
    stats refresh 5s
\end{lstlisting}

\vspace{0.4cm}

\begin{lstlisting}[language=bash, caption={Konfigurasi haproxy-wlc.cfg (Weighted Least Connections)}]
backend web_backends
    balance leastconn
    option httpchk GET /health
    http-check expect status 200
    # Alokasi bobot 1:2:4 dengan algoritma leastconn
    server web1 web-backend-1:8080 check weight 1
    server web2 web-backend-2:8080 check weight 2
    server web3 web-backend-3:8080 check weight 4
\end{lstlisting}

\newpage

% ==============================================================================
% LAMPIRAN 4: KODE SUMBER BACKEND SERVICE GOLANG & DOCKERFILE
% ==============================================================================
\section*{Lampiran 4: Implementasi Layanan Backend Golang \& Berkas Citra Docker}

Berikut adalah implementasi peladen \textit{backend} berbasis bahasa pemrograman Go (\texttt{main.go}) yang mengekspos \textit{endpoint} intensif \texttt{/api/cpu} dan antrean soket \texttt{/api/io}, serta berkas \texttt{Dockerfile} berbasis kompilasi multi-tahap (\textit{multi-stage build}):

\begin{lstlisting}[language=C, caption={Kode Sumber Backend Service (main.go)}]
package main

import (
	"crypto/sha256"
	"encoding/hex"
	"encoding/json"
	"fmt"
	"log"
	"net/http"
	"os"
	"strconv"
	"time"
)

type Response struct {
	Status        string `json:"status"`
	ContainerID   string `json:"container_id"`
	Endpoint      string `json:"endpoint"`
	ExecutionTime string `json:"execution_time_ms"`
	Timestamp     int64  `json:"timestamp"`
	Result        string `json:"result,omitempty"`
}

var containerID string

func init() {
	var err error
	containerID, err = os.Hostname()
	if err != nil {
		containerID = "unknown-container"
	}
}

func healthHandler(w http.ResponseWriter, r *http.Request) {
	w.Header().Set("Content-Type", "application/json")
	w.Header().Set("X-Backend-Server", containerID)
	json.NewEncoder(w).Encode(Response{
		Status:      "healthy",
		ContainerID: containerID,
		Endpoint:    "/health",
		Timestamp:   time.Now().UnixMilli(),
	})
}

// cpuHandler menjalankan hash SHA-256 berulang untuk memicu komputasi vCPU
func cpuHandler(w http.ResponseWriter, r *http.Request) {
	start := time.Now()
	iterations := 10000
	if iterParam := r.URL.Query().Get("iter"); iterParam != "" {
		if val, err := strconv.Atoi(iterParam); err == nil && val > 0 {
			iterations = val
		}
	}

	data := []byte("uniba-madura-skripsi-benchmarking-seed-" + strconv.FormatInt(start.UnixNano(), 10))
	for i := 0; i < iterations; i++ {
		hash := sha256.Sum256(data)
		data = hash[:]
	}

	elapsed := float64(time.Since(start).Microseconds()) / 1000.0

	w.Header().Set("Content-Type", "application/json")
	w.Header().Set("X-Backend-Server", containerID)
	json.NewEncoder(w).Encode(Response{
		Status:        "success",
		ContainerID:   containerID,
		Endpoint:      "/api/cpu",
		ExecutionTime: fmt.Sprintf("%.2f", elapsed),
		Timestamp:     time.Now().UnixMilli(),
		Result:        hex.EncodeToString(data),
	})
}

// ioHandler mensimulasikan latensi operasi basis data (150ms non-blocking)
func ioHandler(w http.ResponseWriter, r *http.Request) {
	start := time.Now()
	delay := 150 * time.Millisecond
	if delayParam := r.URL.Query().Get("delay"); delayParam != "" {
		if ms, err := strconv.Atoi(delayParam); err == nil && ms >= 0 {
			delay = time.Duration(ms) * time.Millisecond
		}
	}

	time.Sleep(delay)
	elapsed := float64(time.Since(start).Microseconds()) / 1000.0

	w.Header().Set("Content-Type", "application/json")
	w.Header().Set("X-Backend-Server", containerID)
	json.NewEncoder(w).Encode(Response{
		Status:        "success",
		ContainerID:   containerID,
		Endpoint:      "/api/io",
		ExecutionTime: fmt.Sprintf("%.2f", elapsed),
		Timestamp:     time.Now().UnixMilli(),
		Result:        fmt.Sprintf("I/O delay of %v completed", delay),
	})
}

func main() {
	port := os.Getenv("PORT")
	if port == "" {
		port = "8080"
	}

	mux := http.NewServeMux()
	mux.HandleFunc("/", healthHandler)
	mux.HandleFunc("/health", healthHandler)
	mux.HandleFunc("/api/cpu", cpuHandler)
	mux.HandleFunc("/api/io", ioHandler)

	server := &http.Server{
		Addr:         ":" + port,
		Handler:      mux,
		ReadTimeout:  15 * time.Second,
		WriteTimeout: 15 * time.Second,
		IdleTimeout:  60 * time.Second,
	}

	log.Printf("[INFO] Backend container %s listening on port %s...", containerID, port)
	if err := server.ListenAndServe(); err != nil && err != http.ErrServerClosed {
		log.Fatalf("[FATAL] Server error: %v", err)
	}
}
\end{lstlisting}

\vspace{0.4cm}

\begin{lstlisting}[language=bash, caption={Berkas Multi-Stage Dockerfile Backend}]
# Tahap 1: Kompilasi Binary Go
FROM golang:1.22-alpine AS builder
WORKDIR /app
COPY go.mod ./
COPY main.go ./
RUN CGO_ENABLED=0 GOOS=linux go build -ldflags="-s -w" -o backend-app main.go

# Tahap 2: Runtime Image Ringan
FROM alpine:3.20
RUN apk --no-cache add ca-certificates tzdata
WORKDIR /root/
COPY --from=builder /app/backend-app .
ENV PORT=8080
EXPOSE 8080
CMD ["./backend-app"]
\end{lstlisting}

\newpage

% ==============================================================================
% LAMPIRAN 5: SKRIP UJI BEBAN GRAFANA K6
% ==============================================================================
\section*{Lampiran 5: Skrip Skenario Pengujian Beban Grafana k6}

Berikut adalah skrip otomasi pengujian beban stres k6 untuk skenario \textit{CPU-bound} (\texttt{cpu-test.js}) dan \textit{I/O-bound} (\texttt{io-test.js}) dengan penambahan pengguna virtual bertahap (100 $\rightarrow$ 500 $\rightarrow$ 1.000 VUs):

\begin{lstlisting}[language=bash, caption={Skrip Skenario CPU-bound (cpu-test.js)}]
import http from 'k6/http';
import { check, sleep } from 'k6';
import { Counter, Rate, Trend } from 'k6/metrics';

export const reqBackend1 = new Counter('backend_1_requests');
export const reqBackend2 = new Counter('backend_2_requests');
export const reqBackend3 = new Counter('backend_3_requests');
export const errorRate = new Rate('custom_error_rate');
export const cpuLatency = new Trend('cpu_endpoint_duration', true);

export const options = {
    stages: [
        { duration: '1m', target: 100 },  // Tahap pemanasan (Warm-up) 100 VUs
        { duration: '3m', target: 500 },  // Beban menengah 500 VUs
        { duration: '4m', target: 1000 }, // Beban puncak (Peak saturation) 1.000 VUs
        { duration: '2m', target: 0 },    // Penurunan beban (Cool-down)
    ],
    thresholds: {
        http_req_failed: ['rate<0.05'], // Toleransi error rate maksimal 5%
        http_req_duration: ['p(95)<3000', 'p(99)<5000'], // SLO p95 < 3s, p99 < 5s
    },
};

const BASE_URL = __ENV.TARGET_URL || 'http://localhost';

export default function () {
    const res = http.get(`${BASE_URL}/api/cpu`, {
        timeout: '10s',
        tags: { endpoint: 'cpu_bound' },
    });

    const isSuccess = check(res, {
        'status is 200': (r) => r.status === 200,
        'has valid result': (r) => r.body && r.body.includes('success'),
    });

    errorRate.add(!isSuccess);
    cpuLatency.add(res.timings.duration);

    const backendServer = res.headers['X-Backend-Server'];
    if (backendServer) {
        if (backendServer.includes('web-backend-1')) {
            reqBackend1.add(1);
        } else if (backendServer.includes('web-backend-2')) {
            reqBackend2.add(1);
        } else if (backendServer.includes('web-backend-3')) {
            reqBackend3.add(1);
        }
    }

    sleep(0.1);
}
\end{lstlisting}

\newpage

% ==============================================================================
% LAMPIRAN 6: OTOMASI EKSEKUSI MATRIKS & AUDIT MAPE
% ==============================================================================
\section*{Lampiran 6: Skrip Otomasi Pengujian Matriks \& Analisis Presisi MAPE}

Berikut adalah skrip \textit{orchestration runner} (\texttt{run-matrix.sh}) untuk mengeksekusi 16 kombinasi eksperimen dan skrip Python (\texttt{parse-logs.py}) untuk menghitung metrik \textit{Mean Absolute Percentage Error} (MAPE):

\begin{lstlisting}[language=bash, caption={Skrip Otomasi 16 Skenario Matriks (run-matrix.sh)}]
#!/usr/bin/env bash
set -e

DIR="$(cd "$(dirname "${BASH_SOURCE[0]}")/.." >/dev/null 2>&1 && pwd)"
cd "$DIR"

mkdir -p logs/nginx logs/haproxy results

run_scenario() {
    local lb="$1"       # nginx / haproxy
    local algo="$2"     # rr / lc / wrr / wlc
    local workload="$3" # cpu / io

    echo ">>> MENJALANKAN SKENARIO: LB=$lb | ALGO=$algo | WORKLOAD=$workload <<<"
    rm -f logs/nginx/access.log logs/haproxy/haproxy.log

    if [ "$lb" == "nginx" ]; then
        export NGINX_CONF="./loadbalancer/nginx/nginx-${algo}.conf"
        docker compose --profile nginx down --volumes --remove-orphans > /dev/null 2>&1 || true
        docker compose --profile nginx up -d --build
        sleep 5
    else
        export HAPROXY_CONF="./loadbalancer/haproxy/haproxy-${algo}.cfg"
        docker compose --profile haproxy down --volumes --remove-orphans > /dev/null 2>&1 || true
        docker compose --profile haproxy up -d --build
        sleep 5
    fi

    local script="benchmark/${workload}-test.js"
    local output_json="results/${lb}_${algo}_${workload}.json"
    k6 run --out json="$output_json" "$script"

    local algo_type="weighted"
    if [ "$algo" == "rr" ] || [ "$algo" == "lc" ]; then
        algo_type="unweighted"
    fi

    if [ "$lb" == "nginx" ]; then
        python3 scripts/parse-logs.py --log logs/nginx/access.log --type nginx --algorithm "$algo_type"
    else
        python3 scripts/parse-logs.py --log logs/haproxy/haproxy.log --type haproxy --algorithm "$algo_type"
    fi
}

run_scenario "$1" "$2" "$3"
\end{lstlisting}

\vspace{0.4cm}

\begin{lstlisting}[language=Python, caption={Skrip Penghitung Deviasi Distribusi Beban MAPE (parse-logs.py)}]
#!/usr/bin/env python3
import sys, re, argparse
from collections import Counter

def parse_nginx_log(log_path):
    counts = Counter()
    pattern = re.compile(r'upstream_addr="([^"]+)"')
    total = 0
    with open(log_path, 'r', encoding='utf-8', errors='ignore') as f:
        for line in f:
            match = pattern.search(line)
            if match:
                addr = match.group(1).split(',')[0].strip()
                counts[addr] += 1
                total += 1
    return counts, total

def calculate_mape(actual_counts, is_weighted=True):
    if is_weighted:
        p_ideal = [100.0 * 1.0 / 7.0, 100.0 * 2.0 / 7.0, 100.0 * 4.0 / 7.0]
    else:
        p_ideal = [100.0 / 3.0, 100.0 / 3.0, 100.0 / 3.0]

    total_reqs = sum(actual_counts.values())
    if total_reqs == 0: return

    nodes_sorted = sorted(actual_counts.keys())
    deviations = []
    for idx, node in enumerate(nodes_sorted):
        vol = actual_counts[node]
        p_act = (vol / total_reqs) * 100.0
        p_id = p_ideal[idx] if idx < len(p_ideal) else (100.0 / len(nodes_sorted))
        dev = abs(p_act - p_id) / p_id * 100.0
        deviations.append(dev)

    mape = sum(deviations) / len(deviations) if deviations else 0.0
    print(f"Mean Absolute Percentage Error (MAPE): {mape:.2f}%")
\end{lstlisting}

\newpage

% ==============================================================================
% LAMPIRAN 7: SKRIP OTOMASI PENGUJIAN MATRIKS DAN TUNING KERNEL LINUX
% ==============================================================================
\section*{Lampiran 7: Skrip Otomasi Pengujian Matriks dan Tuning Kernel Linux}

Berikut adalah skrip otomasi eksekusi matriks pengujian beban dengan isolasi CPU core hibrida (\texttt{run-matrix.sh}) serta konfigurasi optimalisasi soket TCP kernel Linux (\texttt{99-skripsi.conf}):

\begin{lstlisting}[language=bash, caption={Skrip Otomasi Matriks Pengujian Beban dan Isolasi CPU (run-matrix.sh)}]
#!/usr/bin/env bash
# ==============================================================================
# AUTOMATED MATRIX TEST RUNNER (SKRIPSI AHMAD MUQTAFI)
# Lingkungan: Bare-Metal Laptop MSI Thin GF63 12UCX (Ubuntu 24.04 LTS)
# Isolasi CPU: Server pada P-Cores (0-7), k6 pada E-Cores (8-11)
# ==============================================================================
set -e

DIR="$(cd "$(dirname "${BASH_SOURCE[0]}")/.." >/dev/null 2>&1 && pwd)"
cd "$DIR"

mkdir -p logs/nginx logs/haproxy results

# Deteksi ketersediaan core CPU untuk taskset E-Cores (CPU 8-11)
TOTAL_CPUS=$(nproc 2>/dev/null || echo 1)
TASKSET_CMD=""
if [ "$TOTAL_CPUS" -ge 12 ]; then
    TASKSET_CMD="taskset -c 8-11"
    echo "[ISOLASI CPU] Grafana k6 dipin ke E-Cores (CPU 8-11) via taskset."
fi

run_scenario() {
    local lb="$1" algo="$2" workload="$3" quick_flag="$4"
    rm -f logs/nginx/access.log logs/haproxy/haproxy.log

    if [ "$lb" == "nginx" ]; then
        export NGINX_CONF="./loadbalancer/nginx/nginx-${algo}.conf"
        docker compose --profile nginx down --volumes --remove-orphans > /dev/null 2>&1 || true
        docker compose --profile nginx up -d --build
    else
        export HAPROXY_CONF="./loadbalancer/haproxy/haproxy-${algo}.cfg"
        docker compose --profile haproxy down --volumes --remove-orphans > /dev/null 2>&1 || true
        docker compose --profile haproxy up -d --build
    fi
    sleep 5

    local script="benchmark/${workload}-test.js"
    local output_json="results/${lb}_${algo}_${workload}_$(date +%Y%m%d_%H%M%S).json"
    local k6_extra_args=""
    [ "$quick_flag" == "--quick" ] && k6_extra_args="--vus 50 --duration 1m"

    if command -v k6 &> /dev/null; then
        $TASKSET_CMD k6 run $k6_extra_args --out json="$output_json" "$script"
    else
        docker run --rm -i --network="host" --cpuset-cpus="8-11" \
            -v "$DIR/benchmark:/benchmark" grafana/k6 run $k6_extra_args "/$script"
    fi
}
\end{lstlisting}

\vspace{0.4cm}

\begin{lstlisting}[language=bash, caption={Konfigurasi Tuning Parameter Kernel Linux (/etc/sysctl.d/99-skripsi.conf)}]
# Optimalisasi Batas Antrean Koneksi TCP dan Socket Buffer
net.core.somaxconn = 65535
net.ipv4.tcp_max_syn_backlog = 65535
net.ipv4.ip_local_port_range = 1024 65535
net.ipv4.tcp_tw_reuse = 1
net.ipv4.tcp_fin_timeout = 15
fs.file-max = 2097152
\end{lstlisting}


\end{document}
